\section{IV. 주권자와 수범자(SOVEREIGN AND
SUBJECT)}\label{iv.-uxc8fcuxad8cuxc790uxc640-uxc218uxbc94uxc790sovereign-and-subject}

강제적 명령으로서의 법이라는 단순한 모델을 비판하면서, 우리는 지금까지
이 이해방식에 따르면 어떤 사회의 법을 구성하는 것으로 여겨지는
`주권자(sovereign)' 인격 또는 인격들의 존재에 관해서는 아무런 문제도
제기하지 않았다. 실제로, 위협에 의해 뒷받침되는 명령이라는 개념이 법의
다양한 유형들을 설명하기에 충분한지 여부를 논의함에 있어, 우리는
잠정적으로 다음과 같은 가정을 두었다. 즉, 법이 존재하는 모든 사회에는
실제로 하나의 주권자가 존재하며, 그는 복종의 습관(habit of obedience)을
참조하여(by reference to) 다음과 같이 긍정적·부정적으로 특징지어질 수
있다: 사회의 대다수가 그의 명령을 습관적으로 복종하는 반면, 그 자신은
다른 어떤 개인이나 집단의 명령에도 습관적으로 복종하지 않는 인격 또는
인격들의 집단.

이제 우리는 모든 법체계의 기초에 관한 이 일반 이론을 다소 상세히
검토해야 한다. 극단적으로 단순함에도 불구하고, 주권 교설(doctrine of
sovereignty)은 결코 그 중요성이 단순하지 않기 때문이다. 이 이론은, 법이
존재하는 모든 인간 사회에서는, 민주정이든 절대군주정이든 정치적 형태의
다양성 아래에 잠재되어 있는 것으로서, 수범자들이 습관적으로 복종을
행하고 누구에게도 습관적으로 복종하지 않는 주권자 사이의 이 단순한
관계가 궁극적으로 잠재되어 발견된다고 단언한다. 이 이론에 따르면,
주권자와 수범자로 구성된 이러한 수직적 구조는, 법을 보유한 사회에 있어
인간에게 척추가 필수적인 것과 마찬가지로 필수적인 요소이다. 이러한
구조가 현존한다면, 우리는 그 사회를 그 주권자와 함께 하나의 독립된
국가라고 부를 수 있으며, \emphksans{그 사회의} 법(\emph{its} law)에 대해서도
말할 수 있다. 반대로 이러한 구조가 존재하지 않는 경우에는, 이러한 표현들
가운데 어느 것도 적용할 수 없는데, 이는 주권자와 수범자의 관계가 이
이론에 따르면 그러한 표현들의 의미 자체의 일부를 구성하기 때문이다.

이 교설에는 특별히 중요한 두 가지 논점이 있으며, 우리는 이 장의 나머지
부분에서 상세하게 추구될 비판의 방향을 제시하기 위해 여기서 이를
일반적인 용어로 강조하고자 한다. 첫째 논점은 주권자의 법(sovereign's
laws)이 적용되는 자들에게 \textbf{(p.~51)} 요구되는 전부라고 여겨지는
복종의 \emphksans{습관}(a \emph{habit} of obedience)이라는 관념에 관한
것이다. 여기서 우리는 그러한 습관이 대부분의 법체계의 두 가지 두드러진
특징 --- 즉 서로 다른 입법자들이 연속적으로 보유하는 입법 권위(authority
to make law)의 \emphksans{연속성(continuity)}과, 그 법을 만든 자와 그에게
습관적 복종(habitual obedience)을 부여했던 자들이 사망한 이후에도 법이
존속하는 \emphksans{지속성(persistence)} --- 을 설명하기에 충분한지 여부를
검토할 것이다. 둘째 논점은 법 위에 존재하는 주권자의 지위(position)에
관한 것이다. 그는 타인을 위해 법을 창설함으로써 그들에게 법적 책무 또는
`제한(limitations)'을 부과하는 반면, 그 자신은 법적으로 무제한적이며
제한 불가능한(legally unlimited and illimitable) 존재로 간주된다. 여기서
우리는 최고 입법자(supreme lawgiver)의 이러한 법적으로 제한 불가능한
지위가 법의 존재(existence of law)에 필수적인지 여부, 그리고 입법적
권한(legislative power)에 대한 법적 한계(legal limits)의 존재 여부가, 이
이론이 이러한 개념들을 습관과 복종이라는 단순한 용어로 분석하는 방식
속에서 이해될 수 있는지를 검토할 것이다.

\subsection{1. 복종의 습관과 법의
연속성}\label{uxbcf5uxc885uxc758-uxc2b5uxad00uxacfc-uxbc95uxc758-uxc5f0uxc18duxc131}

복종(obedience)이라는 관념은, 엄밀한 검토없이 사용되는 다른 많은
겉보기에는 단순한 관념들과 마찬가지로, 복잡성으로부터 자유롭지 않다.
우리는 이미 지적된 복잡성\footnote{See p.~19 above.}, 즉 `복종'이라는
단어가 종종 위협에 의해 뒷받침된 명령에 대한 단순한 준수(compliance)만이
아니라 권위에 대한 존중을 시사한다는 점은 여기서는 논외로 하겠다.
그럼에도 불구하고, 한 사람이 다른 사람에게 대면하여 단일한 명령(single
order)을 내리는 경우조차도, 특정된 행위(specified act)의
수행(performance)이 복종을 구성하기 위해서는 그 명령의 제시와 그 행위의
수행 사이에 어떠한 정확한 연결이 존재해야 하는지를 명확히 진술하기는
쉽지 않다. 예컨대, 명령을 받은 사람이 어떠한 명령도 없었더라도 바로 그
동일한 일을 확실히 했을 것이라는 사실이, 사실로서 그러한 경우, 어떠한
관련성을 가지는가? 이러한 어려움은 특히 법의 경우에 더욱 첨예하게
나타나는데, 그중 일부는 많은 사람들이 애초에 할 생각조차 하지 않았을
행위를 금지하기 때문이다. 이러한 어려움들이 해결되기 전까지는, 한 나라의
법에 대한 `일반적 복종의 습관(general habit of obedience)'이라는 전체적
관념은 다소 불분명한 상태로 남을 수밖에 없다. 그러나 우리의 현재 목적을
위해서는, `습관'과 `복종'이라는 단어들이 비교적 분명하게 적용될 수
있다고 인정될 법한 매우 단순한 사례를 상정해 볼 수 있을 것이다.

\textbf{(p.~52)} 우리는 다음과 같은 상황을 가정해 보자. 어떤 영토에
인구가 거주하고 있고, 그곳에서는 절대군주인 군주(Rex)가 매우 오랜 기간
동안 통치하고 있다. 그는 일반적 명령을 위협으로 뒷받침하여 사람들을
통제하며, 그들이 그렇지 않았다면 하지 않았을 여러 행위를 하도록
요구하고, 또 그렇지 않았다면 했을 행위를 하지 않도록 요구한다. 통치
초기에는 혼란이 있었으나, 오래전에 사태는 안정되었고, 일반적으로
사람들은 그에게 복종할 것이라고 기대할 수 있다. Rex가 요구하는 내용은
종종 부담스러운 것이며, 복종하지 않고 처벌을 감수하고자 하는 유혹도
상당하기 때문에, 일반적으로 복종이 이루어진다 하더라도 그것이 이 단어의
온전한 의미나 가장 통상적인 의미에서의 `습관(habit)' 또는
`습관적인(habitual)' 복종이라고 보기는 어렵다. 사람들은 실제로 어떤 법에
대한 준수하기 습관을 문자 그대로의 의미에서 획득할 수 있다. 예컨대
영국인에게 도로의 좌측으로 운전하는 것은 그러한 획득된 습관의 전형적인
사례일 것이다. 그러나 세금 납부를 요구하는 법처럼 강한 성향에 반하는
경우에는, 비록 우리가 그것을, 심지어 정기적으로(regular), 결과적으로
준수(compliance)하더라도 그 준수는 습관이 지니는 무반성적이고, 노력이
들지 않으며, 몸에 배어 있는 성격을 갖지 않는다. 그럼에도 불구하고
Rex에게 부여되는 복종은 흔히 이러한 습관적 요소를 결여하고 있을지라도,
다른 중요한 요소들은 지니게 된다. 어떤 사람이 예컨대 아침 식사 때 신문을
읽는 습관을 가지고 있다고 말하는 것은, 그가 상당한 기간 동안 그렇게
해왔고, 또한 그러한 행위를 반복할 개연성이 있다는(likely) 것을 함축한다.
그렇다면 상정된 이 공동체의 대부분의 사람들에 대하여, 초기의 혼란기를
지난 어느 시점 이후에는, 그들이 일반적으로 Rex의 명령에 복종해 왔고
앞으로도 계속 그렇게 할 개연성이 있다는(likely) 점은 참일 것이다.

여기서 주목할 점은, Rex의 통치 하에서의 이 사회적 상황에 대한 설명에
따르면, 복종의 습관(habit of obedience)은 Rex와 각 수범자 간의 개별적
관계라는 점이다. 각 개인은, Rex가 여러 사람에게 내린 명령 중 자신에게
해당하는 것을 정기적으로(regularly) 수행한다. 우리가 ``이 \emphksans{인구
전체(population)}가 그런 습관을 갖고 있다''고 말할 경우, 이는 ``사람들이
토요일 밤마다 주점에 가는 습관이 있다''고 말할 때와 마찬가지로, 단지
대다수 사람들의 습관이 수렴한다(convergent)는 의미일 뿐이다. 즉, 그들
각각이 Rex의 명령에 복종하듯, 각각이 토요일 밤마다 주점을 찾는 것이다.

이 매우 단순한 상황에서는, 공동체가 Rex를 주권자로 성립시키기 위해
요구되는 것은 오직 인구 구성원들의 개인적 복종 행위(personal acts of
obedience)뿐이라는 점에 유의할 필요가 있다. 각 개인은 자기 몫으로 단지
복종하기만 하면 되고, \textbf{(p.~53)} 복종이 정기적으로(regularly)
이루어지기만 한다면, 공동체의 누구도 자기 자신의 복종이나 타인의 Rex에
대한 복종이 어떤 의미에서 옳은지(right), 적절한지(proper), 또는 정당하게
요구되는 것인지(legitimately demanded)에 관하여 어떠한 견해를 가질
필요도, 그것을 표현할 필요도 없다. 분명히 우리가 묘사한 이 사회는 복종의
습관(habit of obedience)이라는 개념에 가능한 한 문자 그대로 적용하기
위해 설정된, 매우 단순한 사회이다. 이는 아마도 실제로는 어디에서도
존재한 적이 없을 만큼 지나치게 단순한 사회일 가능성이 크며, 또한 결코
원시 사회도 아니다. 왜냐하면 원시 사회는 Rex와 같은 절대적
지배자(absolute rulers)를 거의 알지 못하고, 그 구성원들 역시 단지
복종하는 데에만 관심을 두기보다는, 관련된 모든 사람들의 복종이 옳은지에
관하여 뚜렷한 견해를 가지는 것이 보통이기 때문이다. 그럼에도 불구하고,
Rex의 지배하에 있는 이 공동체는 적어도 Rex의 생존 기간 동안에는 법에
의해 지배되는 사회가 지니는 몇 가지 중요한 표지를 분명히 갖추고 있다. 이
공동체는 심지어 일정한 통일성(unity)도 지니고 있어서 `국가(state)'라고
불릴 수도 있다. 이러한 통일성은 구성원들이 복종이 옳은지에 대한 견해를
전혀 공유하지 않더라도, 모두가 동일한 한 사람에게 복종한다는 사실에 의해
구성된다.

이제 성공적인 통치 이후 렉스(Rex)가 사망하고, 그의 아들 Rex 2세가 남아
일반적 명령(general orders)을 발하기 시작했다고 가정해 보자. Rex 1세의
생전 동안 그에게 일반적인 복종의 습관(general habit of obedience)이
존재했다는 사실만으로는, 그것만으로 Rex 2세 역시 습관적으로 복종을 받게
될 가망조차(probable) 말할 수 없다. 따라서 우리가 Rex 1세에 대한 복종의
사실과, 그가 계속 복종을 받을 개연성(likelihood)만을 근거로 삼는다면,
Rex 1세의 마지막 명령에 대해서는 말할 수 있었던 것---즉 그것이
주권자(sovereign)에 의해 내려졌고 따라서 법(law)이었다는 점---을 Rex
2세의 첫 번째 명령에 대해서는 말할 수 없게 된다. 아직 Rex 2세에 대한
확립된 복종의 습관이 존재하지 않기 때문이다. 우리는 이 이론에 따르면,
Rex 2세가 이제 주권자이며 그의 명령이 법이라고 말할 수 있으려면, 그의
부친에게 그랬던 것처럼 Rex 2세에게도 복종이 실제로 부여될 것인지를
기다려 보아야 한다. 출발 시점에서 그를 주권자로 만들어 주는 것은
아무것도 없다. 그의 명령들이 일정 기간 동안 복종을 받아 왔다는 사실을
알게 된 이후에야 비로소, 복종의 습관이 확립되었다고 말할 수 있을 것이다.
그때가 되어서야, 그리고 그때에 이르러서야, 우리는 이후에 내려지는 어떤
명령에 대해서도 그것이 발해지는 즉시, 그리고 실제로 복종되기 이전에도
이미 법이라고 말할 수 있게 된다. 이 단계에 도달하기 전까지는, 어떠한
법도 제정될 수 없는 공백기(interregnum)가 존재하게 된다.

이러한 사태는 물론 가능하며 혼란의 시기에는 때때로 현실화되기도 했다.
그러나 단절(discontinuity)의 위험은 분명하며, 보통은 이를 감수하려 하지
않는다. 그 대신, \textbf{(p.~54)} 절대군주제하에서도 법체계(legal
system)의 특징은 한 입법자에서 다른 입법자로의 이행을 연결하는 규칙들에
의해 입법권(law-making power)의 중단 없는 연속성(uninterrupted
continuity)을 확보하는 데 있다. 이러한 규칙들은 입법자의 자격과 그를
결정하는 방식(mode)을 \emphksans{사전에(in advance)} 규율하며, 이를
지명하거나 일반적 용어로 특정한다. 현대 민주국가에서는 이러한 자격이
매우 복잡하고 구성원이 빈번히 교체되는 입법부의 구성과 관련되지만,
연속성을 위해 요구되는 규칙의 본질은 우리가 상정한 가상의 군주제에
적합한 보다 단순한 형태에서도 확인될 수 있다. 만약 규칙이 장자의 계승을
규정한다면, Rex 2세는 부친의 뒤를 이을 \emphksans{권원(title)}을 가진다. 그는
부친의 사망 시 법을 제정할 \emphksans{권리(right)}를 가지며, 그의 첫 명령이
발해질 때에는 그 개인과 수범자들 사이에 습관적 복종(habitual
obedience)의 관계가 성립할 시간이 아직 없었더라도, 그것들이 이미
법이라고 말할 좋은 이유(good reason)가 있을 수 있다. 실제로 그러한
관계는 끝내 성립하지 않을 수도 있다. 그럼에도 그의 말은 법이 될 수 있다.
왜냐하면 Rex 2세는 첫 명령을 발한 직후 곧바로 사망할 수도 있으며, 그는
복종을 받으며 살 기회를 갖지 못했더라도 법을 제정할 권리를 가졌고, 그의
명령은 법일 수 있기 때문이다.

개별 입법자(individual legislators)의 교체되는 승계를 통해
입법권(law-making power)의 연속성을 설명할 때, `승계 규칙(rule of
succession)', `권원(title)', `승계할 권리(right to succeed)', `법을
제정할 권리(right to make law)'와 같은 표현을 사용하는 것은 자연스럽다.
그러나 분명한 것은, 이러한 표현들과 함께 우리는 복종의 습관(habits of
obedience)에 기초한 일반적 명령(general orders)만으로는 설명될 수 없는
새로운 요소들의 집합을 도입하게 된다는 점이다. 주권 이론(theory of
sovereignty)의 처방을 따라 Rex 1세의 단순한 법적 세계를 구성할 때
사용했던 그러한 요소들로는, 이 새로운 요소들을 설명할 수 없다. 왜냐하면
그 세계에는 규칙이 존재하지 않았고, 따라서 권리나 권원도 존재하지
않았으며, 그러므로 \emphksans{더더군다나(a fortiori)} 승계할 권리나 권원은
존재하지 않았기 때문이다. 그곳에는 단지 Rex 1세가 명령을 발했고, 그의
명령이 습관적으로 복종되었다는 사실만이 있었다. 그의 생존 기간 동안
Rex를 주권자로 구성하고 그의 명령을 법으로 만드는 데에는 그 이상이
필요하지 않았지만, 이것만으로는 그의 후계자의 \emphksans{권리들(rights)}을
설명하기에는 충분하지 않다. 실제로, 복종의 습관이라는 관념은 한 입법자가
다른 입법자를 계승할 때 모든 정상적인 법체계(normal legal system)에서
관찰되는 연속성(continuity)을 두 가지 서로 연관되지만 구별되는 방식으로
설명하지 못한다. 첫째, 한 입법자가 발한 명령에 대한 단순한 복종의
습관만으로는 \textbf{(p.~55)} 새로운 입법자에게 구(舊) 입법자를 계승하여
그의 자리에 서고 명령을 발할 어떤 \emphksans{권리(right)}도 부여할 수 없다.
둘째, 구(舊) 입법자에 대한 습관적 복종만으로는 새로운 입법자의 명령이
복종될 것이라는 가망(probability)을 그 자체로 확보하거나, 그러한
추정(presumption)을 기초지울 수 없다. 만약 승계의 순간에 이러한 권리와
이러한 추정이 존재하려면, 이전 입법자의 통치 기간 동안 사회 어딘가에는
복종의 습관이라는 개념으로는 설명될 수 없는, 그보다 더 복잡한 일반적
사회적 실천(general social practice)이 존재했어야 한다. 즉, 새로운
입법자가 승계할 자격을 가지게 되는 규칙의 수용(acceptance of a rule)이
있었어야 한다.

그렇다면 이러한 보다 더 복합적인 실천(more complex practice)이란
무엇인가? `규칙의 수용(acceptance of a rule)'이란 무엇인가? 여기서
우리는 제1장에서 이미 개괄해 두었던 탐구를 다시 이어가야 한다. 이 질문에
답하기 위해 우리는 한동안 법 규칙의 특수한 경우에서 벗어나야 한다.
습관(habit)과 규칙(rule)은 어떻게 다른가? 사람들이 토요일 밤마다
영화관에 가는 습관이 있다고 말하는 것과, 교회에 들어설 때 남성은 반드시
모자를 벗어야 한다는 규칙이 있다고 말하는 것 사이에는 어떤 차이가
있는가? 우리는 이미 제1장에서 이러한 유형의 규칙을 분석하기 위해 반드시
도입되어야 할 몇 가지 요소들을 언급한 바 있으며, 이제는 그 분석을 한
걸음 더 진전시켜야 할 것이다.

사회적 규칙과 습관 사이에는 확실히 하나의 유사점이 존재한다. 두 경우
모두에서 문제 되는 행위(예: 교회에서 모자를 벗는 행위)는 반드시
전반적으로 일반적이어야 하며, 반드시 항상 불변할 필요는 없다. 즉, 기회가
생기면 대다수 집단 구성원이 그 행위를 반복한다는 점에서 ``그들은 그것을
\emphksans{규칙으로써(as a rule)} 한다''고 말하는 표현은 바로 이 의미를 담고
있다. 하지만 이러한 유사성이 있음에도 불구하고, 양자 사이에는 세 가지
중요한 차이점이 있다.

첫째, 집단이 어떤 \emphksans{습관(habit)}을 갖고 있다고 하기 위해서는, 단지
그들의 행동이 실제로 수렴하면(converge) 충분하다. 즉, 정기적인(regular)
코스에서 벗어나는 일이 있어도 그것이 어떤 비판의 대상이 될 필요는 없다.
그러나 행동의 일반적 수렴 또는 일치만으로는, 그 행동을 요구하는 규칙이
존재한다고 말하기에 충분하지 않다. 규칙이 있는 경우에는, 그로부터의
이탈은 일반적으로 실수(lapse) 또는 과실(fault)로 간주되어 비판의 대상이
되고, 이탈이 예상될 경우 동조를 요구하는 압력이 발생한다. 이러한 비판과
압력의 양상은 규칙의 유형에 따라 다양하다.

둘째로, 그러한 규칙들이 존재하는 경우에는, 실제로 그러한 비판이 이루어질
뿐만 아니라, 그 기준에서의 이탈 자체가 일반적으로 그러한 비판을 제기하기
위한 하나의 \emphksans{좋은 이유(good reason)}로 받아들여진다. 기준에서의
이탈에 대한 비판은 \textbf{(p.~56)} 이와 같은 의미에서
정당하거나(legitimate) 정당화된(justified) 것으로 간주되며, 이탈이
예상될 때 그 기준에의 준수를 요구하는 요구(demands) 또한 마찬가지로
정당한 것으로 간주된다. 더 나아가, 완고한 소수의 상습적 위반자(hardened
offenders)를 제외하면, 이러한 비판과 요구는 그것을 제기하는 사람들뿐만
아니라 그것을 받는 사람들에 의해서도 일반적으로 정당한 것으로, 또는 좋은
이유에 근거한 것으로 받아들여진다. 이러한 여러 방식으로 집단 구성원들 중
몇 명이, 그리고 얼마나 자주 또 얼마나 오랫동안, 행태의 정기적인
방식(regular mode of behaviour)을 비판의 기준(standard of criticism)으로
취급해야 비로소 그 집단이 하나의 규칙(rule)을 가지고 있다고 말할 수
있는지는 명확히 정해진 문제가 아니다. 이는 사람이 대머리라고 불릴 수
있으면서도 머리카락을 몇 가닥이나 가질 수 있는지에 관한 문제와
마찬가지로, 우리를 지나치게 괴롭힐 필요는 없다. 우리가 기억해야 할 것은,
어떤 집단이 특정한 규칙을 가지고 있다는 진술은, 그 규칙을 위반할 뿐만
아니라 그 규칙을 자기 자신이나 타인에 대해 기준으로 보기를 거부하는
소수가 존재하는 것과도 양립 가능하다는 점이다.

사회적 규칙(social rules)을 습관(habits)과 구별하는 세 번째 특징은 이미
앞에서 말한 내용 속에 암묵적으로 포함되어 있지만, 법리학에서 매우
중요함에도 불구하고 자주 간과되거나 오인되어 온 것이므로 여기서 이를 좀
더 상세히 설명하고자 한다. 이 특징을 우리는 이 책 전반에 걸쳐 규칙의
\emphksans{내적 측면(internal aspect)}이라고 부를 것이다. 어떤 습관이 사회
집단 내에서 일반화되어 있을 때, 이러한 일반성은 그 집단 구성원 대다수의
관찰 가능한 행동에 관한 하나의 사실에 불과하다. 그러한 습관이 존재하기
위해서는, 집단의 구성원들 중 누구도 그 일반적인 행동을 어떤 방식으로든
의식할 필요가 없고, 나아가 문제의 행동이 일반적이라는 사실을 알 필요조차
없다. 더더욱 그 행동을 가르치거나 유지하려고 노력하거나 의도할 필요도
없다. 각자가 다른 사람들도 실제로 그렇게 행동하고 있다는 점에서 자신
또한 그렇게 행동하기만 하면 충분하다. 이에 반해, 사회적 규칙이
존재하려면 적어도 일부 구성원들은 문제의 행동을 집단 전체가 따라야 할
일반적 기준(general standard)으로 간주해야 한다. 사회적 규칙은 사회적
습관과 공유하는 외적 측면(external aspect), 즉 관찰자가 기록할 수 있는
정기적이고 획일적인 행동(regular uniform behaviour)이라는 측면에 더하여,
이러한 `내적' 측면을 함께 가진다.

이러한 규칙의 내적 측면은 어떤 게임의 규칙에서 쉽게 예시할 수 있다. 체스
플레이어들은 단지 외부 관찰자가 ``퀸은 저렇게 움직이는구나''라고 기록할
수 있을 정도로 비슷한 방식으로 퀸을 움직이는 습관만 갖고 있는 것이
아니다. 이에 더해, \textbf{(p.~57)} 그들은 이 행동 양식에 대해
성찰적이고 비판적인 태도를 취한다. 즉, 그것을 체스를 두는 사람 모두가
따라야 할 기준으로 본다. 각자는 단지 퀸을 일정한 방식으로 움직이는 것에
그치지 않고, 모두가 그렇게 해야 한다고 생각한다. 이러한 생각은, 다른
사람이 규칙을 어기거나 어길 가능성이 있을 때, 그를 비판하고 동조를
요구하는 방식으로 표현되며, 다른 사람이 자신에게 비판이나 요구를 할 때
이를 정당한 것으로 받아들이는 방식에서도 드러난다. 이러한 비판, 요구,
수용의 표현에는 다양한 규범적 언어(normative language)가 사용된다.
예컨대 ``퀸을 그렇게 움직이면 안 됐어(I/You ought not to have moved the
Queen like that)'', ``그렇게 해야 해(I/You must do that)'', ``그건
맞아(That is right)'', ``그건 틀렸어(That is wrong)'' 등이 이에
해당한다.

규칙의 내적 측면(internal aspect)은 종종 외적으로 관찰 가능한 물리적
행동(physical behaviour)에 대비되는 단순한 `감정(feelings)'의 문제로
오도되어 이해된다. 물론 규칙이 어떤 사회 집단에 의해 일반적으로
수용되고, 사회적 비판과 순응을 요구하는 압력에 의해 일반적으로 지지되는
경우, 개인들은 제약(restriction)이나 강제(compulsion)의 경험과 유사한
심리적 경험을 흔히 가질 수 있다. 사람들이 자신이 특정한 방식으로
행동하도록 `구속되어 있다고 느낀다(feel bound)'고 말할 때, 실제로는
이러한 경험들을 가리키고 있을 수도 있다. 그러나 이러한 감정들은 `구속력
있는(binding)' 규칙의 존재를 위해 필요조건도 아니고 충분조건도 아니다.
사람들이 특정한 규칙을 수용하면서도 그러한 강제의 감정을 전혀 경험하지
않는다고 말하는 데에는 아무런 모순도 없다. 필요한 것은, 일정한 행동
패턴(patterns of behaviour)을 공동의 기준(common standard)으로 간주하는
비판적 성찰적 태도가 존재하는 것이며, 이러한 태도가 비판(자기비판을
포함하여), 순응 요구(demands for conformity), 그리고 그러한 비판과
요구가 정당화될 수 있다는 인정(acknowledgements) 속에서 드러나는 것이다.
이 모든 것은 `\textasciitilde 해야 한다(ought)', `\textasciitilde 해야만
한다(must)', `\textasciitilde 하는 것이 당연하다(should)',
`옳다(right)'와 `그르다(wrong)'라는 규범적 용어(normative
terminology)에서 그 특유한 표현을 발견한다.

이것들이 바로 사회적 규칙(social rules)을 단순한 집단적 습관(mere group
habits)과 구별하는 핵심적 특징들이며, 이러한 점들을 염두에 두고 우리는
다시 법으로 돌아갈 수 있다. 우리는 우리의 사회 집단이 교회에서 모자를
벗는 것에 관한 규칙처럼 특정한 유형의 행위를 기준(standard)으로 만드는
규칙들뿐만 아니라, 일정한 사람의 말---구두로 표현되었든 문서로
표현되었든---을 참조함으로써(by reference to) 보다 간접적인 방식으로
행위의 기준을 식별하도록 하는 규칙을 가지고 있다고 가정할 수 있다. 가장
단순한 형태에서 이러한 \textbf{(p.~58)} 규칙은, Rex가 (아마도 일정한
형식적 방식으로) 지정하는 어떠한 행위들이 행해져야 한다는 내용일 것이다.
이는 우리가 처음에 Rex에 대한 단순한 복종의 습관이라는 관점에서 묘사했던
상황을 변형시킨다. 왜냐하면 이러한 규칙이 수용되는 경우, Rex는 단지
사실상 무엇을 해야 하는지를 지정하는 데 그치지 않고, 그렇게 할
\emphksans{권리}(the \emph{right})를 가지게 되며, 그의 명령에 대한 일반적인
복종이 존재할 뿐 아니라, 그에게 복종하는 것이 \emphksans{옳다(right)}(it is
\emph{right} to obey him)는 점 또한 일반적으로 수용되기 때문이다. Rex는
사실상 입법할 \emphksans{권위(authority)}를 가진 입법자(legislator)가 되며,
즉 집단의 삶 속에 새로운 행위 기준(standards of behaviour)을 도입할
권한을 갖게 된다. 그리고 이제 우리가 `명령'이 아니라 기준(standards)을
다루고 있는 이상, 그가 자신의 입법에 의해 구속되어서는 안 될 이유는 전혀
없다.

이러한 입법적 권위(legislative authority)의 기초를 이루는 사회적
실천들(social practices)은 교회에서 머리를 벗는 행위에 관한 규칙과 같은
단순하고 직접적인 행위 규칙(simple direct rules of conduct)의 기초를
이루는 실천들과, 모든 본질적인 점들에 있어, 동일하며, 우리는 이제 이를
단순한 관습 규칙(mere customary rules)으로 구별할 수 있고, 또한 일반적
습관(general habits)과는 동일한 방식으로 구별된다. 이제 Rex의
\emphksans{말(word)}은 행위의 기준이 되어, 그가 지정한 행위로부터의 이탈은
비판의 대상이 될 수 있으며; 그의 말은 일반적으로 비판과 준수
요구(demands for compliance)를 정당화하는 근거로 참조되고 수용될 것이다.

이러한 규칙들이 입법적 권위의 연속성(continuity of legislative
authority)을 어떻게 설명하는지를 이해하기 위해서는, 어떤 경우에는 새로운
입법자가 실제로 입법을 시작하기도 전에, 그가 하나의 \emphksans{계층(class)}
또는 계보(line)에 속한 사람으로서 차례가 되었을 때 입법할 권리(right)를
가진다는 점을 부여하는 확고히 확립된 규칙이 존재함이 이미 분명할 수
있다는 점에 주목하기만 하면 된다. 예컨대 Rex 1세의 생존 기간 동안, 집단
전체에 의해 일반적으로 수용되는 바에 따르면, 그 말(word)이 복종의 대상이
되는 사람은 개별 인물 Rex 1세에 한정되는 것이 아니라, 일정한 방식으로
자격을 갖춘 사람, 예컨대 특정 조상의 직계(line) 중 생존한 최연장자로서
그 시점에 자격을 갖춘 바로 그 사람이라는 것이다. Rex 1세는 단지 특정한
시점에 그러한 자격을 갖춘 특정 인물에 불과하다. 이러한 규칙은 Rex
1세에게 복종하는 습관과는 달리, 현재의 실제 입법자뿐만 아니라 장차
가능하게 등장할 미래의 입법자들까지 지시한다는 점에서 미래를 향해 있다.

이러한 규칙의 수용(acceptance), 그리고 따라서 그 존재(existence)는 Rex
1세의 생존 기간 동안 부분적으로는 그에 대한 복종을 통해 드러나지만, 또한
일반 규칙(general rule)에 따른 그의 자격(qualification)에 의해 복종이
그가 권리를 가진 어떤 것이라는 점에 대한 인정(acknowledgements)을
통해서도 드러난다. \textbf{(p.~59)} 어떤 집단에 의해 특정 시점에 수용된
규칙의 범위(scope)가 이와 같은 방식으로 일반적으로 입법자 공직(office of
legislator)의 후임자를 향해 열려 있을 수 있다는 바로 그 이유 때문에, 그
규칙의 수용은 후임자가 실제로 입법을 시작하기 이전부터 입법할 권리(right
to legislate)를 가진다는 법적 진술(statement of law)에 대해서도, 그리고
그가 전임자와 동일한 복종을 받을 개연성이 있다는(likely) 사실
진술(statement of fact)에 대해서도 \emphksans{둘 모두(both)}에 대한
근거(grounds)를 제공한다.

물론, 어떤 사회가 어느 한 시점에서 규칙을 수용한다고 해서 그 규칙의
지속적인 존재가 \emphksans{보장되는(guarantee)} 것은 아니다. 혁명이 일어날 수
있으며, 그 사회가 해당 규칙을 더 이상 수용하지 않게 될 수도 있다. 이는
한 입법자, Rex 1세의 생존 기간 중에 발생할 수도 있고, 새로운 입법자 Rex
2세로의 이행 시점에서 발생할 수도 있으며, 만약 그러한 일이 발생한다면
Rex 1세는 입법할 권리(right to legislate)를 상실하게 되거나, Rex 2세는
그 권리를 취득하지 못하게 된다. 상황이 불분명해질 수도 있다는 점은
사실이다. 즉, 단순한 봉기(insurrection)나 기존 규칙의 일시적
중단(temporary interruption)에 직면해 있는 것인지, 아니면 그것의
전면적이고 실효적인 포기(full-scale effective abandonment)에 직면해 있는
것인지가 명확하지 않은 중간적이고 혼란스러운 단계가 존재할 수 있다.
그러나 원칙적으로 이 문제는 분명하다. 새로운 입법자가 입법할 권리를
가진다는 진술(statement)은, 그가 그 권리를 가지게 하는 규칙이 사회 집단
내에 존재함을 전제한다. 만약 현재 그에게 자격을 부여하는 규칙이,
동일하게 그의 전임자에게도 자격을 부여했던 것으로서, 전임자의 생존 기간
동안 수용되었음이 분명하다면, 반대 증거가 없는 한, 그 규칙은 포기되지
않았으며 여전히 존재하고 있다고 가정되어야 한다. 유사한 연속성은
경기(game)에서도 관찰된다. 즉, 득점 기록자는 지난 이닝(last innings)
이후 경기 규칙이 변경되었다는 증거가 없는 한, 새로운 타자가 통상적인
방식으로 평가되어 만들어낸 출루를 그에게 귀속시킨다.

Rex 1세와 Rex 2세의 단순한 법적 세계를 고찰하는 것만으로도, 대부분의
법체계를 특징짓는 입법 권위(legislative authority)의 연속성이 한 규칙의
수용을 구성하는 사회적 실천의 한 형태에 의존하며, 우리가 이미 지적한
방식들에서 단순한 습관적 복종의 사실들과는 구별된다는 점을 보여주기에
충분할 것이다. 논증은 다음과 같이 요약될 수 있다. Rex와 같이 그의 일반적
명령이 습관적으로 복종되는 사람이 입법자(legislator)라 불릴 수 있고 그의
명령들이 법이라 불릴 수 있다는 점을 우리가 인정하더라도(concede), 그러한
입법자들의 연속(succession)에 속한 각자에 대한 복종의 습관만으로는
후임자(successor)가 계승할 \emphksans{권리(right)}와 그에 따른 입법적 권한의
연속성을 설명하기에 충분하지 않다. 첫째, \textbf{(p.~60)} 습관은
`규범적이지(normative)' 않다. 즉, 그것들은 누구에게도 권리(rights)나
권위(authority)를 부여할 수 없다. 둘째, 한 개인에 대한 복종의 습관은,
수용된 규칙(accepted rules)과는 달리, 현재의 입법자뿐만 아니라 장차
연속적으로 등장할 입법자들의 계보(class or line of future successive
legislators)을 참조하거나(refer to), 그들에 대한 복종을
개연적으로(likely) 만들 수 없다. 따라서 한 입법자에 대한 습관적 복종이
존재한다는 사실은, 그의 후임자가 입법할 권리를 가진다는
진술(statement)에 대해서도, 그가 실제로 복종을 받을 개연성이
있다는(likely) 사실 진술(factual statement)에 대해서도, 아무런
근거(grounds)를 제공하지 않는다.

그러나 이 지점에서, 우리가 뒤의 장에서 충분히 전개할 중요한 논점 하나를
주목하지 않을 수 없다. 이는 오스틴의 이론이 지니는 강점 가운데 하나를
이룬다. 수용된 규칙(accepted rules)과 습관(habits) 사이의 본질적 차이를
드러내기 위해 우리는 매우 단순한 형태의 사회를 가정해 왔다.
주권(sovereignty)의 이 측면을 떠나기 전에, 입법할 권위를 부여하는
규칙(rule conferring authority to legislate)의 수용에 대한 우리의 설명이
어느 정도까지 현대적 국가(modern state)에 얼마나 적용될 수 있을지
검토해야 한다. 우리의 단순한 사회를 언급하면서 우리는 마치 대부분의 일반
시민들이 법에 단지 복종할 뿐만 아니라, 법 제정자의 계승을 정당화하는
규칙을 이해하고 수용하고 있는 것처럼 말했다. 단순한 사회에서는 이것이
사실일 수 있다. 그러나 현대 국가에서는, 아무리 법을 잘
지키는(law-abiding) 사람들이라 하더라도, 끊임없이 변화하는 입법권한이
부여된 사람들의 집단의 자격(qualifications)을 구체화하는 규칙들에 대해
인구 대다수가 명확한 인식을 가지고 있다고 생각하는 것은 터무니없을
것이다. 어떤 소규모 부족의 구성원들이 그들의 연속적인 추장에게 권위를
부여하는 규칙을 수용하는 것과 동일한 방식으로, 대중이 이러한 규칙들을
`수용한다(accept)'고 말하는 것은, 보통의 시민들이 실제로는 갖고 있지
않을지도 모르는 헌법적 사안에 대한 이해를 그들의 머릿속에 집어넣는 것을
함의하게 된다. 우리는 이러한 이해를 체계의 공무담당자들 또는
전문가들에게만 요구하면 된다. 즉, 법이 무엇인지를 결정할
책임(responsibility)을 부여받은 법원(courts)과, 일반 시민이 법이
무엇인지 알고자 할 때 자문하는 변호사들(lawyers)이 그것이다.

단순한 부족 사회와 현대 국가 사이의 이러한 차이들은 주목할 가치가 있다.
그렇다면, 연속적으로 교체되는 입법자들의 변화 속에서도 유지되는 의회의
여왕(Queen in Parliament)의 입법적 권위의 연속성을, 일반적으로 수용된
어떤 근본 규칙 또는 규칙들에 기초한 것으로 우리는 어떤 의미에서 이해해야
하는가? 분명히 여기서의 일반적 수용(general acceptance)은
\textbf{(p.~61)} 복합적인 현상으로서, 공무담당자와 일반 시민 사이에 어떤
의미에서 분할되어 있으며, 이들은 서로 다른 방식으로 이에 기여함으로써
법체계의 \emphksans{존재(existence)}에 기여한다. 체계의 공무담당자들은 입법
권위를 부여하는 이러한 근본 규칙들을 명시적으로 인정(acknowledge
explicitly)한다고 말할 수 있다. 입법자들은 자신들에게 권한을
부여하는(empower) 규칙들에 따라 법을 제정할 때 그렇게 하고, 법원은
그러한 자격을 갖춘 자들에 의해 제정된 법을 자신들이 적용할 법으로
식별할(identify) 때 그렇게 하며, 전문가들은 그렇게 제정된 법을 참조하여
일반 시민들의 행위를 향도할(guide) 때 그렇게 한다. 일반 시민은 이러한
공적 작용(official operations)의 결과에 대한 묵인(acquiescence)을 통해
자신의 수용을 주로 표출한다. 그는 이러한 방식으로 제정되고 식별된 법을
지키고(keeps the law), 그 법에 의해 부여된 청구(claims)를 제기하며
권한(powers)을 행사한다. 그러나 그는 그 법의 기원(origin)이나
제정자(makers)에 대해서는 거의 알지 못할 수도 있다. 어떤 이들은 법에
대해 그것이 단지 `법(the law)'이라는 점 이상을 전혀 알지 못할 수도 있다.
법은 일반 시민들이 하고 싶어 하는 일들을 금지하며(forbids), 그들은 이를
위반할 경우 경찰 공무원에 의해 체포되고 판사에 의해 징역형에 처해질 수
있다는 사실을 알고 있다. 위협에 의해 뒷받침된 명령에 대한 습관적 복종이
법체계의 기초라는 점을 고집하는 이 교설의 강점은, 우리가 법체계의
존재라고 부르는 이 복합적 현상의 상대적으로 수동적인 측면을 현실적인
용어로 사유하도록 강제한다는 데 있다. 이 학설의 약점은, 체계의 공직자나
전문가들의 법 제정(law-making), 법 식별(law-identifying), 법
적용(law-applying) 작용에서 주로---비록 배타적으로는
아니지만---드러나는, 다른 상대적으로 능동적인 측면을 흐리게 하거나
왜곡한다는 데 있다. 우리가 이 복합적인 사회적 현상을 실제로 그러한
것으로 이해하려면, 이 두 측면을 모두 염두에 두어야 한다.

\subsection{2. 법의 지속성}\label{uxbc95uxc758-uxc9c0uxc18duxc131}

1944년, 한 여성은 영국에서 점을 치는 행위로 1735년 제정된
「주술법(Witchcraft Act)」을 위반한 혐의로 기소되어 유죄 판결을
받았다.\footnote{\emph{R.} v. \emph{Duncan}{[}1944{]} 1 KB 713.} 이는
매우 익숙한 법적 현상의 하나를 보여주는 다소 인상적인 사례에 불과하다.
즉, 수세기 전에 제정된 법률이 오늘날에도 여전히 법일 수 있다는 것이다.
그러나 이 현상이 아무리 익숙하다고 하더라도, 이러한 방식의 법의
지속성(persistence)은, \textbf{(p.~62)} 법을 습관적으로 복종받는 한
인물(person)이 내린 명령으로 파악하는 단순한 도식으로는 이해 가능하게
만들 수 없다. 우리는 여기서 사실상 방금 살펴본 입법 권위의
연속성(continuity) 문제의 역(converse)에 직면해 있다. 앞의 경우에는,
복종의 습관이라는 단순한 도식에 기초하여, 입법자 지위(office)의
후임자(successor)가 제정한 최초의 법이, 그가 개인적으로 아직 습관적
복종을 받기 이전임에도 \emphksans{이미(already)} 법이라고 말할 수 있는지가
문제였다. 여기서의 문제는 이와 반대로, 오래전에 사망한 선행 입법자가
제정한 법이, 그에게 습관적으로 복종한다고 말할 수 없는 사회에서 어떻게
\emphksans{여전히(still)} 법일 수 있는가 하는 것이다. 첫 번째 경우와
마찬가지로, 우리가 입법자의 생존 기간(lifetime of the legislator)으로
시야를 한정한다면, 이 단순한 도식에는 아무런 어려움도 발생하지 않는다.
실제로 그것은 왜 「주술법」이 영국에서는 법이었지만 프랑스에서는 법이
아니었는지를 매우 훌륭하게 설명하는 것처럼 보인다. 설령 그 법의 용어들이
프랑스에서 점을 치는 프랑스 시민에게까지 확장되었다 하더라도
마찬가지이다. 물론 불운하게도 영국 법원에 회부된 프랑스인들에게는 그
법이 적용될 수 있었을 것이다. 이 단순한 설명에 따르면, 영국에서는 이
법을 제정한 자들에 대한 복종의 습관이 존재했으나, 프랑스에서는 그러한
습관이 존재하지 않았다는 것이다. 따라서 이 법은 영국에서는 법이었지만
프랑스에서는 법이 아니었던 것이다.

그러나 우리는 법을 그 제정자들의 생존 기간으로만 한정하여 볼 수는 없다.
우리가 설명해야 할 특징은 바로 법이 그 제정자들과 그들에게 습관적으로
복종하던 사람들까지 사라진 이후에도 끈질기게 존속할 수 있는
능력(capacity)에 있기 때문이다. 「주술법」이 동시대의 프랑스인들에게는
법이 아니었다면, 왜 그것이 오늘날 우리에게는 여전히 법인 것인가? 분명히,
어떻게 언어를 늘여 써도, 20세기의 영국인들인 우리가 지금 George 2세와
그의 의회(Parliament)에 습관적으로 복종한다고 말할 수는 없다. 이 점에서
오늘날의 영국인들과 당시의 프랑스인들은 동일하다. 즉, 어느 쪽도 이 법의
제정자에게 습관적으로 복종하지 않았거나 복종하고 있지 않다. 「주술법」이
이 재위 기간 동안 제정된 유일한 법률이라고 하더라도, 그것은 오늘날에도
여전히 영국에서 법일 것이다. 이 `왜 여전히 법인가?'라는 문제에 대한
해답은, 원칙적으로, 우리가 앞서 제기한 `왜 이미 법인가?'라는 첫 번째
문제에 대한 해답과 동일하다. 그것은 주권적 인물(sovereign person)에 대한
복종의 습관이라는 지나치게 단순한 개념을, 사회에서 행동의 기준(standard
of behaviour)을 구성할 말(word)을 발할 자들의 계급(class) 또는 계보(line
of persons), 다시 말해 입법할 \emphksans{권리(right)}를 가진 자들을 규정하는,
현재 수용된 근본 규칙들(currently accepted fundamental rules)이라는
개념으로 대체하는 것을 포함한다. 그러한 규칙은 비록 현재(now) 존재해야
하지만, 어떤 의미에서는 시간 초월적(timeless)인 참조(reference)를 가질
수 있다. 즉, 그것은 \textbf{(p.~63)} 장차의 입법자(future legislator)의
입법 작용(legislative operation)을 지시하며 앞으로 나아갈 뿐만 아니라,
과거의 입법자(past one)의 작용을 지시하며 뒤를 돌아볼 수도 있다.

Rex 왕조라는 단순한 틀로 제시하면, 상황은 다음과 같다. Rex 1세, 2세,
3세로 이어지는 입법자들의 계보에 속한 각자는, 직계(direct line)에 있는
생존한 최연장 후손에게 입법할 권리를 부여하는 동일한 일반 규칙(general
rule)에 따라 자격을 갖출 수 있다. 개별 통치자(individual ruler)가
사망하면 그의 입법 작업(legislative work)은 계속 존속한다. 이는 그것이,
각 입법자가 생존해 있던 때마다 사회의 연속적인 세대(successive
generations)가 존중해 온 일반 규칙이라는 기초 위에 놓여 있기 때문이다.
단순한 사례에서 Rex 1세, 2세, 3세는 모두 동일한 일반 규칙에 따라, 입법을
통해 행동 기준(standards of behaviour)을 도입할 자격을 가진다. 대부분의
법체계에서는 사정이 이처럼 단순하지 않다. 과거의 입법을 법으로 인정하는
현재 수용된 규칙(presently accepted rule)이, 동시대의 입법에 관한 규칙과
서로 다를 수도 있기 때문이다. 그러나 그 기초가 되는 규칙(underlying
rule)이 현재 수용되고 있다는 점을 전제로 한다면, 법의 지속성(persistence
of laws)은, 구성원이 변경된 팀들 사이의 토너먼트 첫 라운드에서 내려진
심판의 판정(decision)이, 세 번째 라운드에서 그 자리를 대신한 심판의
판정과 마찬가지로 최종 결과(final result)에 동일한 관련성(relevance)을
갖는다는 사실만큼이나 신비로운 것이 아니다. 그럼에도 불구하고,
신비롭지는 않다 하더라도, 과거와 미래, 그리고 현재의 입법자들의
명령에까지 권위를 부여하는 수용된 규칙의 관념은, 현재의 입법자에 대한
복종의 습관이라는 관념보다 분명히 더 복잡하고 정교하다. 그렇다면 이러한
복잡성을 제거하고, 위협에 의해 뒷받침된 명령이라는 단순한 개념을 어떤
기지 있는 확장(ingenious extension)을 통해 발전시켜, 법의 지속성이
결국에는 현재의 주권자에 대한 습관적 복종이라는 더 단순한 사실들 위에
기초하고 있음을 보여줄 수는 없는가?

이를 달성하려는 한 가지 기지 있는 시도가 제시된 바 있다. 홉스(Hobbes)는,
여기서 벤담(Bentham)과 오스틴에 의해 반복되듯이, ``입법자(legislator)란
법이 처음 제정되었을 때 그 법에 권위(authority)를 부여한 자가 아니라, 그
법이 지금도 법으로 존속하도록 권위를 부여하는 자이다''라고
말하였다.\footnote{\emph{Leviathan}, chap.~xxvi.} 규칙의 개념을 배제하고
보다 단순한 습관의 관념을 택한다면, 입법자의 `권위(authority)'가 그의
`권한(power)'과 구별되어 무엇을 의미할 수 있는지는 즉각적으로 명확하지
않다. 그러나 이 인용문에 표현된 일반적 \textbf{(p.~64)} 논증은 분명하다.
그것은 「주술법」과 같은 법의 원천(source)이나 기원(origin)이 역사적
사실로서 과거의 주권자(past sovereign)의 입법 작용에 있었음에도
불구하고, 20세기 잉글랜드에서 그 법이 현재 법으로서의 지위(status as
law)를 가지는 것은 현재의 주권자에 의해 법으로 승인(recognition)되기
때문이라는 주장이다. 이러한 인정은, 현재 살아있는 입법자들에 의해 제정된
법률의 경우와 같은 \emphksans{명시적(explicit)} 명령의 형태를 취하지 않고,
주권자의 의지의 \emphksans{묵시적(tacit)} 표현의 형태를 취한다. 이는,
주권자가 그럴 수 있음에도 불구하고, 오래전에 제정된 그 법률을 자신의
대리인들(agents)-법원과 경우에 따라서는 행정부-에 의해 집행되는 것을
방해하지 않는다는 사실로 구성된다.

이는 물론, 어느 누구에 의해서도 어느 시점에서도 명령된 것처럼 보이지
않는 일정한 관습적 규칙(customary rules)의 법적 지위를 설명하기 위해
원용되었던, 이미 검토된 묵시적 명령에 관한 동일한 이론이다. 제III장에서
우리가 이 이론에 대해 제기한 비판들은, 과거의 입법(past legislation)이
법으로서 계속 승인되는 것을 설명하는 데 이 이론이 사용될 때 더욱
분명하게 적용된다. 법원이 합당하지 않은(unreasonable) 관습 규칙을 배제할
수 있는 광범위한 재량(wide discretion)을 부여받고 있기 때문에, 특정
사건에서 법원이 실제로 관습 규칙을 적용하기 전까지는 그것이 법으로서의
지위를 갖지 않는다는 견해에는 어느 정도의 그럴듯함(plausibility)이 있을
수 있다. 그러나 과거의 `주권자'에 의해 제정된 법률이, 특정 사건에서
법원에 의해 실제로 적용되고, 현재의 주권자의 묵인 아래 집행되기 전까지는
법이 아니라는 견해에는 거의 그럴듯함이 없다. 만약 이 이론이 옳다면,
법원은 그것이 이미 법이기 때문에 집행하는 것이 아니라는 결론이 따라야
한다. 그러나 이러한 결론은, 현재의 입법자가 과거의 제정(enactments)을
폐지(repeal)할 수 있었음에도 그 권한(power)을 행사하지 않았다는
사실로부터 도출하기에는 터무니없이 부당한(absurd) 추론이다. 빅토리아
시대 법률들과 오늘날 의회의 여왕(Queen in Parliament)에 의해 제정된
법률들은 분명히 현재의 영국에서 정확히 동일한 법적 지위(legal status)를
가진다. 양자는 모두, 그것들이 적용될 사건들이 법원에 제기되기 이전부터
이미 법이다. 그리고 그러한 사건들이 발생했을 때, 법원은 빅토리아 시대
법률과 현대 법률을 모두, 그것들이 이미 법이기 때문에 적용한다. 어느
경우에서도 이들 법률은 법원에 의해 적용된 이후에야 비로소 법이 되는 것이
아니며, 양 경우 모두에서 그 법으로서의 지위는, 이 법률들이 현재 수용된
규칙들에 따라 권위를 가지는(authoritative) 자들에 의해 제정되었다는
사실에 기인하며, 그 사람들이 생존해 있는지 사망했는지의 여부와는
무관하다.

\textbf{(p.~65)} 과거의 제정법이 현재 법으로서의 지위를 법원에 의한
적용에 대해 현 입법부가 묵인하고 있기 때문에 갖게 된다는 이론의
비일관성(incoherence)은, 오늘날의 법원들이 폐지되지(repealed) 않은
빅토리아 시대 법률을 여전히 법으로 취급하는 반면, Edward 7세 치하에서
폐지된 법률을 더 이상 법이 아닌 것으로 구별해야 하는 이유를 설명하지
못한다는 점에서 가장 분명하게 드러난다. 분명히 이러한 구별을 할 때,
법원들은(그리고 그와 함께 이 체계를 이해하는 변호사나 일반 시민들도)
과거와 현재의 입법 작용을 모두 포괄하는, 무엇이 법으로 간주되어야
하는지를 규정하는 근본 규칙(fundamental rule) 또는 규칙들(rules)을
기준(criterion)으로 사용한다: 그들은 두 법률 사이의 구별을, 현재의
주권자가 한쪽은 묵시적으로 사령(즉, 집행을 허용)하고 다른 쪽은 그렇게
하지 않았다는 사실에 대한 앎(knowledge)에 근거시키지 않는다.

다시 말해, 우리가 배척한 이 이론에서 유일하게 미덕이라고 할 만한 것은
현실주의적 요소를 흐릿한 형태로 상기시켜 준다(a realistic reminder)는
점뿐인 듯하다. 이 경우 그 상기(reminder)란, 해당 체계의 공무담당자들,
특히 무엇보다도 법원이, 특정한 입법 작용이 \emphksans{과거이든 현재이든}
권위를 가진다는 규칙을 수용하지 않는다면, 법으로서의 지위에 필수적인
어떤 것이 결여될 것이라는 점을 상기시키는 것이다. 그러나 이러한 단조로운
유형의 현실주의는, 그 주요 특징들이 뒤에서 상세히 논의되는\footnote{See
  pp.~136--47 below.} 이른바 법적 현실주의(Legal Realism)라는 이론으로
과도하게 확대되어서는 안 된다. 이 이론은 일부 버전에서, 어떤 제정법도
법원에 의해 실제로 적용되기 전까지는 \emphksans{어떠한(no)}
제정법(statutes)도 법이 \emphksans{아니라}고 주장한다. 어떤 법률이 법이 되기
위해서는 법원이 특정한 입법 작용이 법을 만든다는 규칙을 수용해야 한다는
진실(truth)과 법원이 특정 사건에서 그것을 적용하기 전까지는 아무것도
법이 아니라는 오도된 이론(misleading theory) 사이에는, 법을 이해하는 데
있어 결정적으로 중요한 차이가 있다. 물론 법적 현실주의 이론의 일부
버전들은, 우리가 비판해 온 법의 지속성에 대한 거짓 설명보다도 한참
넘어선다. 왜냐하면 그것들은 과거의 주권자에 의해 제정되었든
\emphksans{현재의} 주권자에 의해 제정되었든, 법원이 실제로 적용하기 전에는
어떠한 제정법도 법의 지위(status of law)를 가질 수 없다고 전면적으로
부정하기 때문이다. 그러나 완전한 현실주의 이론(full Realist
theory)에까지 이르지는 않으면서도, 과거의 주권자와 구별하여 현재의
주권자가 제정한 제정법은 법원에 의해 적용되기 이전에도 법이라고 인정하는
한편, 법의 지속성을 설명하려는 이론은 \textbf{(p.~66)} 양쪽 세계의
최악을 취하는 것이며, 분명히 터무니없이 부조리하다. 이러한 중간적 입장은
옹호될 수 없다. 왜냐하면 현재의 주권자가 제정한 법률의 법적 지위와, 이전
주권자에 의해 제정되었으나 폐지되지 않은 법률의 법적 지위를 구별할
아무런 근거도 없기 때문이다. 둘 다(보통의 변호사들이 인정하듯이) 법원이
오늘날의 특정 사건에 적용하기 이전부터 법이거나, 아니면 완전한 현실주의
이론이 주장하듯이 둘 다 법이 아니어야 한다.

\subsection{3. 입법권에 대한 법적
한계}\label{uxc785uxbc95uxad8cuxc5d0-uxb300uxd55c-uxbc95uxc801-uxd55cuxacc4}

주권 교설에서 수범자의 일반적 복종의 습관은 주권자에게서의 그러한 습관의
부재(absence)를, 그것의 상보적 짝(complement)으로, 가지게 된다. 그는
수범자들을 위해 법을 만들며, 어떠한 법에도 속하지 않는 지위(position
outside any law)에서 법을 만든다. 그의 법-생성 권한(law-creating
power)에는 법적 한계(legal limits)이 존재하지 않으며, 또 존재할 수도
없다. 주권자의 법적으로 무제한적인 권한은 정의상(by definition) 그에게
귀속된다는 점을 이해하는 것이 중요하다: 이 이론은 단순히 다음을 단언할
뿐인데, 즉 입법자에게 법적 제한이 존재할 수 있는 경우는 오직 그 입법자가
자신이 습관적으로 복종하는 다른 입법자의 명령 아래에 있을 때뿐이라는
것과, 그 경우 그는 더 이상 주권자가 아니라는 것이다. 그가 주권자라면
그는 어떤 다른 입법자에게도 복종하지 않으며, 따라서 그의 입법 권력에는
법적 제한이 존재할 수 없다. 물론 이 이론의 중요성은, 사실에 대해
아무것도 말해 주지 않는 이러한 정의(definitions)와 그로부터 도출되는
단순하고 필연적인 귀결들(simple necessary consequences)에 있는 것은
아니다. 그 중요성은, 법이 존재하는 모든 사회에는 이러한
속성(attributes)을 지닌 주권자가 존재한다는 다음과 같은 주장(claim)에
있다: 모든 법적 권력이 제한되어 있으며, 어떠한 개인이나 집단도
주권자에게 귀속되는 어떠한 법에도 속하지 않는 지위를 차지하고 있지 않은
것처럼 보이게 하는 법적 또는 정치적 형식 배후를 우리는 들여다볼 필요가
있을지도 모른다. 그러나 우리가 집요하게 추적한다면, 이 이론이 주장하듯이
그러한 형식들 배후에 실제로 존재하는 현실(reality)을 발견하게 될 것이다.

우리는 이 이론을 실제로 그것이 주장하는 바보다 더 약하거나 더 강한
주장을 하는 것으로 오해해서는 안 된다. 이 이론은 법적 한계을 받지 않는
주권자가 존재하는 \emphksans{어떤(some)} 사회가 존재한다는 점만을 진술하는
것이 아니라, 법의 존재는 어디에서나 그러한 주권자의 존재를 함의한다는
점을 주장한다. 한편 이 이론은 주권자의 권한에 아무런 한계도 없다고
주장하는 것이 아니라, 오직 그 권한에 \emphksans{법적} 한계(\emph{legal}
limits)가 없다고 주장할 뿐이다. 따라서 주권자는 입법 권한을 행사함에
있어 실제로는 대중적 입장에 따라 자제할 수도 있다. \textbf{(p.~67)} 이는
이를 멸시할 경우의 결과에 대한 두려움 때문일 수도 있고, 혹은 자신이
도덕적으로 이를 존중하도록 도덕적으로 구속된다고 스스로 생각하고 있기
때문일 수도 있다. 매우 다양한 요인들이 이러한 점에서 그에게 영향을 미칠
수 있으며, 만약 대중의 반란에 대한 두려움이나 도덕적 확신(moral
conviction) 때문에 그가 그렇지 않았다면 했을 방식으로 입법하지 않게
된다면, 그는 실제로 이러한 요인들을 자신의 권한에 대한 `한계'라고
생각하고 말할 수도 있다. 그러나 그것들은 법적 한계(legal limits)는
아니다. 그는 그러한 입법을 삼가야 할 법적 책무(legal duty)를 지고 있지
않으며, 법원(law courts)은 자신들 앞에 주권자의 법이 존재하는지를
판단함에 있어, 대중적 입장이나 도덕성의 요구로부터의 분기가 그 법을
법으로서 인정받지 못하게 만든다는 주장을, 주권자가 그렇게 하라고 명령한
경우가 아니라면, 받아들이지 않을 것이다.

이 이론이 법에 대한 일반적 설명(general account of law)으로서 지니는
매력은 분명하다. 이 이론은 두 가지 주요한 질문에 대해 만족스러울 만큼
단순한 형태로 답을 제시해 주는 것처럼 보인다. 습관적 복종을 받지만
누구에게도 복종을 돌려주지 않는 주권자를 발견하게 되면, 우리는 두 가지
일을 할 수 있다. 첫째, 우리는 그의 일반적 명령 속에서 특정 사회의 법을
식별할 수 있고, 그 사회 구성원들의 삶을 또한 규율하는 다른 많은
규칙(rules), 원칙(principles), 또는 기준(standards)---도덕적이거나
단순히 관습적인(moral or merely customary) 것들과---을 그것으로부터
구별할 수 있다. 둘째, 법의 영역(area of law) 안에서, 우리가 독립적인
법체계(independent legal system)에 직면해 있는지, 아니면 어떤 더
광범위한 체계(wider system)의 종속적인 부분(subordinate part)에 불과한
것에 직면해 있는지를 판단할 수 있다.

일반적으로, 단일하고 연속적인 입법적 존재자(single continuing
legislative entity)로서 이해되는 의회의 여왕(Queen in Parliament)은 이
이론의 요건(requirements)을 충족시키며, 의회의 주권성은 바로 그러한 충족
사실에 존재한다고 주장된다. 이 믿음의 정확성 여부가 어떠하든(그 일부
측면은 제VI장에서 나중에 검토한다), 우리는 이 이론이 요구하는 바를 Rex
1세의 상상속의 단순화된 세계 속에서 충분히 일관되게 재현할 수 있다. 보다
복잡한 현대 국가의 사례를 검토하기에 앞서 이를 수행하는 것은 교육적으로
유익할 것인데, 이 방식으로 이 이론의 전체 함의(full implications)가 가장
잘 드러나기 때문이다. 제1절에서 복종의 습관 개념에 대해 제기된 비판들을
수용하기 위해, 우리는 상황을 습관이 아니라 규칙의 관점에서 구상할 수
있다. 이러한 발판 아래에서, 우리는 법원(courts), 공무담당자, 시민에 의해
일반적으로 수용된 규칙이 존재하는 사회를 상정할 것이다. 그 규칙에
따르면, Rex가 무엇이든 하라고 명령할 때마다 그의 말(word)은 집단을 위한
행동의 기준(standard of behaviour)을 구성한다. 이러한 명령들 가운데에서,
Rex가 `공적' 지위(`official' status)를 부여하고자 하지 않는
\textbf{(p.~68)} `사적' 바람(`private' wishes)의 표현과, 그러한 지위를
부여하고자 하는 표현을 구별하기 위해, 보조적 규칙(ancillary rules)들이
채택될 수도 있다. 이 규칙들은 군주(monarch)가 `군주로서의 성격으로(in
the character of a monarch)' 입법할 때 사용해야 할 특별한 방식(style)을
규정하되, 아내나 정부(情婦)에게 사적 명령을 내릴 때에는 이를 사용하지
않도록 명시할 것이다. 입법의 방식과 형식(manner and form)에 관한 이러한
규칙들은 그 목적을 달성하기 위해서는 진지하게 받아들여져야 하며, 때로는
Rex에게 불편(inconvenience)을 초래할 수도 있다. 그럼에도 불구하고,
우리가 그것들을 법적 규칙(legal rules)으로 분류할 수는 있겠지만, 그의
입법 권한(legislative power)에 대한 `한계(limits)'로 간주할 필요는 없다.
요구된 형식(required form)을 따르기만 한다면, 그의 바람(wishes)을
실현하기 위해 입법할 수 없는 주제는 존재하지 않기 때문이다. 그의 입법
권력의 `형식(form)'은 법에 의해 한계지어질 수 있을지 모르나, 그
`범위(area)'는 법에 의해 한계지어지지 않는다.

법에 대한 일반 이론(general theory of law)으로서 이 이론에 제기될 수
있는 반론은, 이 상상된 사회에서의 Rex와 같은 어떠한 법적 제한에도
종속되지 않는 주권자의 존재가 법의 존재를 위한 필요조건(necessary
condition)이나 전제(presupposition)가 아니라는 점이다. 이를 입증하기
위해 우리는 논쟁의 여지가 있거나 이의를 제기할 만한 법의 유형들을 원용할
필요가 없다. 따라서 우리의 논증은, 입법부가 결여되어 있다는 이유만으로
법이라는 지위(title of law)를 부정하고자 하는 이들도 있는 관습법 체계나
국제법에서 도출되는 것이 아니다. 이러한 사례들에 호소할 필요는 전혀
없다. 왜냐하면 법적으로 무제한적인 주권자라는 개념은, 법이 존재한다는
점에 대해 누구도 의문을 제기하지 않을 많은 현대 국가들에서의 법의
성격(character of law)을 잘못 대변하기 때문이다. 이러한 국가들에는
입법부들이 존재하지만, 체계 내에서의 최고 입법 권한이 무제한과는 거리가
먼 경우도 종종 있다. 성문헌법(written constitution)은, 입법의 형식과
방식을 규정함으로써---이를 우리는 제한으로 보지 않아도 된다고 인정할 수
있지만---입법부의 권능(competence)을 제한할 뿐만 아니라, 입법 권능의
범위에서 특정 사안들을 전면적으로 배제함으로써, 실체적 제한을 부과할
수도 있다.

다시 한 번, 현대 국가의 복잡한 사례를 검토하기에 앞서, Rex가 최고
입법자인 단순화된 세계에서 그의 입법 권한에 대한 `법적 제한'이 실제로
무엇을 의미하는지, 그리고 왜 그것이 완전히 일관된 생각인지를 살펴보는
것이 유용하다.

Rex의 단순화된 사회에서는, (성문헌법에 구현되어 있든 그렇지 않든 간에)
Rex의 어떠한 법도 토착 주민을 영토(territory)에서 배제하거나 재판 없는
구금을 규정한다면 유효하지 않으며, 이러한 규정에 반하는 어떠한 제정도
\textbf{(p.~69)} 무효(void)로서 간주되고 모두에 의해 그렇게 취급된다는
규칙이 수용되어 있을 수도 있다. 이러한 경우 Rex의 입법할 권한은 분명히
법적(legal)이라고 할 수밖에 없는 제한을 받게 될 것이다. 설령 우리가
이러한 근본적인 헌법 규칙(fundamental constitutional rule)을
`\emphksans{하나의} 법(\emph{a} law)'이라고 부르기를 꺼린다 하더라도 말이다.
그가 종종 자신의 성향(inclinations)에 반하여서까지도 자제하게 할 수 있는
대중적 입장이나 대중의 도덕적 확신(popular moral convictions)을 무시하는
경우와는 달리, 이러한 특정한 제약들(specific restrictions)을 무시하는
것은 그의 입법을 무효로 만들 것이다. 따라서 법원(courts)은, 입법자의
권한 행사에 대한 다른 단지 도덕적 또는 \emphksans{사실상의(de facto)}
한계들과는 달리, 이러한 한계들에 관해서는 직접적인 관련을 갖게 될
것이다. 그럼에도, 이러한 법적 제한들에도 불구하고, 그 범위(scope) 내에서
이루어진 Rex의 제정(enactments)은 분명히 법(laws)이며, 그의 사회에는
독립적인 법체계가 존재한다.

우리는 이러한 상상의 단순한 사례를 좀 더 자세히 들여다볼 필요가 있다.
그래야만 이러한 유형의 법적 제한이 정확히 무엇인지 분명해지기 때문이다.
우리는 종종 Rex의 입장을 다음과 같이 표현할 수 있다: ``그는 재판 없이
구금을 부여하는 법을 제정\emphksans{할 수 없다(cannot)}.'' 이때의
``\textasciitilde 할 수 없다(cannot)''는 표현은, 어떤 사람이 어떤 것을
하지 않아야 할 법적 책무(legal duty)나 의무(obligation)를 지고 있다는
의미와는 다르다. 후자의 의미에서 ``할 수 없다(cannot)''가 사용되는
경우는 예를 들어 ``자전거를 인도에서 탈 수 없다(You cannot ride a
bicycle on the pavement.)''고 말하는 경우이다. 최고 입법자의 권한을
실질적으로 제한하는 헌법은, 입법자에게 특정 방식의 입법을 시도하지
말라는 법적 책무를 부과하지는 않는다(또는 반드시 부과할 필요는 없다). 그
대신, 그런 입법은 무효가 된다고 규정할 수 있다. 이런 경우 입법자에게
부과되는 것은 법적 책무(duties)가 아니라, 법적 무능력(disabilities)이다.
따라서 여기서 `제한(limits)'이란 것은 \emphksans{책무}의 현존(the presence of
\emph{duty})가 아니라 법적 권한의 부재(the absence of legal power)를
의미한다.

Rex의 입법 권한에 대한 이러한 제약들은 헌법적(constitutional)이라고 불릴
수 있을 것이다. 그러나 그것들은 법원이 관여하지 않는 단순한
관행(conventions)이나 도덕적 사안(moral matters)이 아니다. 그것들은
입법할 권위를 부여하는 규칙(rule conferring authority to legislate)의
일부를 이루며, 법원에게 중대한 관련성을 가진다. 왜냐하면 법원은 그러한
규칙을 자신들 앞에 제시되는 입법으로 의도된 제정들(purported legislative
enactments)의 유효성(validity)을 판단하는 기준(criterion)으로 사용하기
때문이다. 그럼에도 불구하고, 이러한 제한들이 법적이며 단순히
도덕적이거나 관습적인 것이 아님에도 불구하고, 그 존재나 부재는 Rex가
다른 사람들에게 습관적으로 복종하는지 여부라는 용어로는 표현될 수 없다.
Rex는 그러한 제약들의 적용을 받을 수 있으며, 그것들을 회피하려는 시도를
\textbf{(p.~70)} 전혀 하지 않을 수도 있다; 그럼에도 불구하고, 그가
습관적으로 복종하는 대상은 아무도 없을 수 있다. 그는 단지 유효한
법(valid law)을 제정하기 위한 조건들(conditions)을 충족시킬 뿐이다. 혹은
그는 그 제약들과 불일치하는 명령(orders inconsistent with them)을
발함으로써 그 제약들을 회피하려고 시도할 수도 있다. 그러나 그가 그렇게
하더라도 그는 누구에게도 불복종한 것이 아니며, 상위 입법자의 법을
위반하거나 법적 책무를 침해한 것도 아니다. 그는 분명히 유효한 법을
제정하는 데 실패했을 뿐이며(그는 법을 어기는 것이 아니라), 그 반대의
경우도 마찬가지이다. 즉, Rex에게 입법 자격을 부여하는 헌법적 규칙 안에
Rex의 입법할 권위에 대한 법적 제약이 전혀 존재하지 않는다면, 그가 인접한
영토의 왕인 Tyrannus의 명령에 습관적으로 복종한다는 사실은, Rex의
제정(enactments)을 법으로서의 지위(status as law)에서 박탈하지도 않으며,
그것들이 Tyrannus가 최고 권위를 가지는 단일 체계(single system)의 종속적
부분(subordinate parts)임을 보여주지도 않는다.

앞서 논의한 명백한 고찰들은, 비록 단순한 주권 이론에 의해 자주 가려져
있지만, 법체계의 기초를 이해하는 데 있어 결정적으로 중요한 여러 사항들을
확립해 준다. 이를 다음과 같이 요약할 수 있다: 첫째, 입법 권한에 대한
법적 제한은 입법자에게 어떤 상위 입법자의 명령에 복종하라는 책무를
부과하는 것이 아니라, 입법 자격을 부여하는 규칙에 포함된
무능력(disabilities)의 형태로 존재한다.

둘째, 어떤 법률 제정이 유효한 법인지 입증하기 위해, 그것이 어떤 `주권자'
혹은 `무제한적' 입법자에 의해 명시적 또는 묵시적으로 제정된 것임을
추적할 필요는 없다. 여기서 `무제한적'이라는 말은, 그의 입법 권한이
법적으로 제한받지 않는다는 의미이거나, 그가 다른 누구에게도 습관적으로
복종하지 않는다는 의미일 수 있다. 그러나 우리가 실제로 해야 할 일은,
해당 제정 행위가 현존하는 규칙에 따라 입법 자격을 갖춘 입법자에 의해
이루어진 것이며, 그 규칙에 어떤 제한이 없다거나, 적어도 해당
제정행위에는 적용되지 않는 제한만이 존재함을 보여주는 것이다.

셋째, 어떤 체계가 독립된 법체계인지 판별하기 위해, 그 최고 입법자가
법적으로 무제한적이거나, 타인에게 습관적으로 복종하지 않는다는 점을
보여줄 필요는 없다. 필요한 것은 단지, 그 입법자를 자격화하는 규칙들이,
다른 영토에 대해서도 권한을 가진 자들에게 우월한 권위를 부여하지
않는다는 사실을 보여주는 것이다. 반대로, 그가 외부 권위에 복종하지
않는다고 해서, 곧바로 그가 자기 영토 내에서 무제한적 권한을 가진다는
뜻은 아니다.

넷째, 법적으로 무제한적인 입법 권한과, \textbf{(p.~71)} 제한은 있지만 그
체계 내에서 최고 권한을 갖는 입법자는 구분되어야 한다. Rex는 헌법에 의해
일정한 제한을 받는다 하더라도, 그의 입법이 다른 모든 입법을 폐지할 수
있는 의미에서 그 나라 법에 의해 인식되는 최고 입법 권한을 가질 수 있다.

다섯째, 그리고 마지막으로, 입법자의 입법 권한에 대한 제한 규칙의 존재
여부는 핵심적인 반면, 그 입법자가 타인에게 습관적으로 복종하는 습관이
있는지 여부는, 기껏해야 간접적인 증거로서만 일정한 의미를 갖는다. 만약
입법자가 다른 사람에게 습관적으로 복종하지 않는다는 것이 사실이라면,
이는 때때로 그 입법자의 권한이 다른 이의 권위에 헌법적 혹은 법적으로
종속되지 않을 수 있다는 점에 대한 약간의, 그러나 결코 결정적인 것은
아닌, 증거로 기능할 수 있다. 마찬가지로, 그 입법자가 실제로 어떤 다른
사람에게 습관적으로 복종한다는 사실 역시, 그의 입법 권한이 어떤 규칙
아래에서 다른 사람의 권위에 종속되어 있음을 나타내는 일정한 증거가 될 수
있다.

\subsection{4. 입법부 배후의
주권자}\label{uxc785uxbc95uxbd80-uxbc30uxd6c4uxc758-uxc8fcuxad8cuxc790}

현대 세계에는, 통상적으로 그 체계 내에서 최고 입법부로 간주되는 기관이
입법 권한의 행사에 있어 법적 제한을 받는 법체계들이 다수 존재한다.
그럼에도 불구하고, 법률가와 법이론가 모두가 동의하듯이, 그러한 입법부가
자신에게 부여된 제한된 권한의 범위(scope) 내에서의 제정(enactments)은
분명히 법(law)이다. 이러한 경우들에서, 만약 우리가 법이 존재하는 곳마다
법적 제한을 받을 수 없는 주권자(sovereign incapable of legal
limitation)가 존재한다는 이론을 유지하고자 한다면, 우리는 법적으로
제한된 입법부 배후에 있는 그러한 주권자를 찾아야 한다. 그가 실제로
거기에 존재하는지 여부가, 이제 우리가 검토해야 할 문제이다.

우리는 잠시 동안, 모든 법체계가 어떤 형태로든---비록 반드시 성문헌법에
의해서일 필요는 없지만---입법자들의 자격(qualification)과 입법의 `방식과
형식'에 관해 마련해야 하는 규정들을 논외로 둘 수 있다. 이러한 규정들은
입법 권력의 범위에 대한 법적 제한이라기보다는, 입법기관의
동일성(identity)과 그것이 입법하기 위해 무엇을 해야 하는지를
특정(specifications)하는 것으로 생각될 수 있다. 그러나 실제로는,
남아프리카(South Africa)의 경험이 보여주듯이,\footnote{See \emph{Harris}
  v. \emph{Dönges}{[}1952{]} 1 TLR 1245.}, \textbf{(p.~72)} 입법의
`방식과 형식'에 관한 단순한 규정이나 입법기관의 정의(definitions)를
`실질적(substantial)' 제한과 만족스럽게 구별해 주는 일반적
기준(criteria)을 제시하기는 어렵다. 그럼에도 불구하고 실질적 제한의
명확한 사례들은, 미합중국이나 오스트레일리아의 연방헌법에서 찾아볼 수
있다. 이들 헌법에서는 중앙정부와 구성 주(states) 사이의 권한
분할(division of powers), 그리고 또한 일정한 개인적 권리(individual
rights)들이 통상의 입법 절차에 의해 변경될 수 없다. 이러한 경우, 주
입법부나 연방 입법부 중 어느 쪽이든, 연방적 권한 분할에 반하거나 그것을
변경하려는 것으로 주장되는 제정(enactment), 또는 이러한 방식으로
보호되는 개인적 권리들과 양립할 수 없는 제정은 \emphksans{권한을 초과한(ultra
vires)} 것으로 취급될 수 있으며, 헌법 규정과 충돌하는 범위 내에서 법원에
의해 법적으로 유효하지 않다(legally invalid)고 선언된다. 이러한 입법
권력에 대한 법적 제한들 가운데 가장 유명한 것은 미합중국 헌법의 수정
제5조(Fifth Amendment)이다. 이 조항은, 다른 내용들과 함께, 어떠한 사람도
적법절차(due process of law) 없이 `생명, 자유 또는 재산(life liberty or
property)'을 박탈당해서는 안 된다고 규정한다. 그리고 의회(Congress)의
법률(statutes)이 이러한 규정들 또는 헌법이 그 입법 권력에 부과한 다른
제한들과 충돌하는 것으로 판명될 경우, 법원에 의해 유효하지 않다고
선언되어(declared invalid) 왔다.

물론 헌법의 규정들을 입법부의 작용으로부터 보호하기 위한 다양한 장치들이
존재한다. 어떤 경우들에서는, 예컨대 스위스의 경우처럼, 연방 구성 주의
권리와 개인의 권리에 관한 일부 규정들이 형식상으로는 강행적임에도
불구하고, `단지 정치적(merely political)'이거나 권고적인(hortatory)
것으로 취급된다. 이러한 경우 법원에는 연방 입법부의 제정을
`심사하여(review)', 비록 그것이 입법부의 작용의 적절한 권한 범위에 관한
헌법 규정과 명백히 충돌하더라도 이를 유효하지 않다(invalid)고 선언할
관할권(jurisdiction)이 부여되지 않는다.\footnote{See Art. 113 of the
  Constitution of Switzerland.} 미합중국 헌법의 특정 규정들 역시 `정치적
질문(political questions)'을 제기하는 것으로 판시되어 왔으며, 사건이 이
범주에 속하는 경우 법원은 해당 법률이 헌법을 위반하는지 여부를 심리하지
않는다.

최고 입법부의 통상적 작용(normal operations)에 대한 법적 제한이 헌법에
의해 부과되는 경우, 이러한 제한들 자체는 \textbf{(p.~73)} 특정한 형태의
법적 변경(legal change)으로부터 면책(immune)될 수도 있고 그렇지 않을
수도 있다. 이는 헌법이 그 개정(amendment)을 위해 어떠한 규정을 두고
있는지의 성격에 달려 있다. 대부분의 헌법은, 통상의 입법부와 구별되는
기관(body)에 의해 행사되거나, 또는 통상의 입법부 구성원들이 특별한
절차(special procedure)를 사용하여 행사하는, 광범위한 개정 권한(amending
power)을 포함하고 있다. 미합중국 헌법 제5조가, 각 주의 입법부 4분의 3
또는 각 주의 4분의 3으로 소집된 회의(conventions)에 의해
비준되는(ratified) 개정을 규정하고 있는 것은 전자의 유형에 속하는 개정
권한의 한 사례이며, 1909년 남아프리카법(South Africa Act of 1909)
제152조(s. 152)에 규정된 개정 절차는 후자의 유형에 속하는 사례이다.
그러나 모든 헌법이 개정 권한을 포함하고 있는 것은 아니며, 또한 개정
권한이 존재하는 경우에도, 입법부에 한계를 부과하는 헌법의 특정 규정들이
그 개정 권한의 범위 밖에 두어지는 경우가 있다. 이 경우 개정 권한 자체가
제한되는 것이다. 이러한 점은(비록 그중 일부 제한은 더 이상 실질적
중요성을 가지지 않지만) 미합중국 헌법에서도 관찰될 수 있다. 왜냐하면
제5조는, `1808년 이전에 이루어진 어떠한 개정도 제1조 제9절의 제1항과
제4항을 어떠한 방식으로도 변경할 수 없으며, 어떠한 주도 그 동의 없이는
상원(Senate)에서의 동등한 참정권(equal suffrage)을 박탈당하지 않는다'고
규정하고 있기 때문이다.

입법부가 남아프리카의 경우처럼, 입법부 구성원들이 특별한 절차를
작동시킴으로써 제거할 수 있는 제한을 받는 경우에는, 그 입법부를 해당
이론이 요구하는 바와 같은 법적 제한을 받을 수 없는 주권자와 동일시할 수
있다는 주장이 가능하다. 이 이론에 있어 곤란한 사례는, 미합중국의
경우처럼 입법부에 대한 제약들(restrictions)이 특별한 기관에 위임된 개정
권한의 행사를 통해서만 제거될 수 있는 경우이거나, 혹은 그러한 제약들이
모두 개정 권한의 범위 밖에 있는 경우이다.

이러한 사례들을 일관되게 설명할 수 있다는 이 이론의 주장을 검토함에
있어, 종종 간과되지만 반드시 상기해야 할 점은, 오스틴 자신이 이 이론을
전개함에 있어 영국에서조차 주권자를 입법부와 동일시하지는
\emphksans{않았다(not)}는 점이다. 이는 의회의 여왕(Queen in Parliament)이
통상적으로 받아들여지는 교설에 따르면 그 입법 권한에 대해 법적
제한으로부터 자유롭고, 따라서 `경성(rigid)' 헌법에 의해 제한되는
의회(Congress)나 다른 입법부들과 대비되는 의미에서 \textbf{(p.~74)}
`주권적 입법부(sovereign legislature)'가 무엇인지를 보여주는
전형(paradigm)으로 자주 인용됨에도 불구하고 그러하였다. 그럼에도
불구하고 오스틴의 견해에 따르면, 어떤 민주정(democracy)에서도 선출된
대표자(elected representatives)가 주권적 기구(sovereign body)를
구성하거나 그 일부를 이루는 것은 아니며, 주권을 구성하는 것은
선거권자들(electors)이다. 따라서 잉글랜드에서는 ``엄밀하게 말하자면
하원의원(members of the commons house)은 자신들을 선출하고 임명한 집단을
위한 단순한 수탁자(trustees)에 불과하며, 결과적으로 주권은 항상
국왕(King), 귀족원(Peers), 그리고 하원의 선거권자 집단(electoral body of
the commons)에 귀속된다''고 하였다.\footnote{Austin, \emph{Province of
  Jurisprudence Determined}, Lecture VI, pp.~230--1.} 마찬가지로 그는
미합중국에서도 각 주의 주권과, ``연방 결합(Federal Union)으로부터 발생한
더 큰 국가의 주권 역시 주 정부들(states' governments)이 하나의 집합적
기구(aggregate body)를 형성하는 데에 귀속되는데, 여기서 주 정부란 통상의
입법부(ordinary legislature)가 아니라, 통상의 입법부를 임명하는 시민들의
집단(body of citizens)을 의미한다''고 보았다.\footnote{Ibid., p.~251.}

이러한 관점에서 보면, 통상의 입법부(ordinary legislature)가 법적
제한으로부터 자유로운 법체계와, 입법부가 그러한 제한을 받는 법체계
사이의 차이는, 주권적 유권자 집단(sovereign electorate)이 자신의 주권적
권한(sovereign powers)을 행사하는 방식의 차이에 불과한 것으로 보인다. 이
이론에 따르면 영국에서는, 유권자 집단이 주권에 참여하는 유일한 직접적
행사는 의회(Parliament)에 앉을 대표자를 선출하고 그들에게 자신의 주권적
권력을 위임하는 데에 있다. 이러한 위임은, 위임된 권력이 남용되지 않도록
신임(trust)이 부여되어 있다는 점에서는 의미를 가지지만, 그러한 신임은
오직 도덕적 제재(moral sanctions)의 문제일 뿐이며, 입법 권력에 대한 법적
제한과 달리 법원이 관여하는 사안은 아니다. 이에 반해, 통상의 입법부가
법적으로 제한되는 모든 민주정에서와 마찬가지로 미합중국에서는, 선거인
집단(electoral body)이 자신의 주권적 권한 행사를 대표자 선출에만
국한하지 않고, 그 대표자들에게 법적 제약(legal restrictions)을 부과해
왔다. 이 경우 유권자 집단은, 헌법적 제약(constitutional restrictions)을
준수할 법적 의무를 지는 통상의 입법부에 비해 우월한, `통상을 넘어서고
이면에 감추어진 입법부(``extraordinary and ulterior legislature'')'로
간주될 수 있으며, 충돌이 발생하는 경우 법원은 통상의 입법부의
조항들(Acts)을 유효하지 않다(invalid)고 선언할 것이다. 따라서 이 경우,
이 이론이 요구하는 바와 같이 모든 법적 제한으로부터 자유로운 주권자는
바로 유권자 집단(electorate)에 존재하게 된다.

\textbf{(p.~75)} 이 이론의 이러한 추가적 전개 국면에서는, 주권자에 대한
초기의 단순한 개념이 급진적인 변형까지는 아니더라도 일정한 정교화를
거쳤다는 점이 분명하다. `사회 구성원의 대다수가 습관적으로 복종하는
대상인 인격 또는 인격들(the person or persons to whom the bulk of the
society are in the habit of obedience)'로서의 주권자에 대한 기술은, 이
장 제1절에서 우리가 보여주었듯이, Rex가 절대 군주였고 입법자로서의
승계(succession)에 관한 어떠한 규정도 존재하지 않았던 가장 단순한 형태의
사회에 대해서는 거의 문자 그대로 적용될 수 있었다. 그러나 그러한 규정이
마련된 경우에는, 현대 법체계의 매우 두드러진 특징인 입법적 권위의
연속성(continuity)은 복종의 습관이라는 단순한 용어로는 표현될 수 없었고,
실제로 입법을 수행하고 복종을 받기 이전에 이미 후임자(successor)가
입법할 권리를 가진다는 점을 규정하는 수용된 규칙(accepted rule)이라는
개념을 필요로 하였다. 그러나 민주국가의 유권자 집단(electorate)과
주권자를 동일시하는 현재의 설명은, 핵심어인 `복종의 습관'과 `인격 또는
인격들'에, 단순한 사례에 적용되었을 때와는 전혀 다른 의미를 부여하지
않는 한, 전혀 그럴듯함을 갖지 못한다. 그리고 이러한 의미는 수용된
규칙이라는 개념이 은밀하게 도입되지 않는 한 명확해질 수 없다. 복종의
습관과 명령이라는 단순한 도식만으로는 이를 감당할 수 없다.

이 점은 여러 방식으로 입증될 수 있지만, 가장 분명하게 드러나는 예는
다음과 같은 경우이다. 유권자에서 유아와 정신적 결함자를 제외한 모든
구성원이 포함되어 있는 민주주의 사회에서는, 유권자 자체가 인구의
`대다수(bulk)'를 구성한다. 또는, 모든 성인이 투표권을 가진 단순한
사회집단을 상상해 보자. 이러한 경우 유권자를 주권자로 간주하고, 원래
이론의 단순한 정의를 적용하려 하면, 우리는 결국 사회의 `대다수'가
스스로에게 습관적으로 복종한다고 말하게 된다. 따라서 ``법적으로 제한을
받지 않는 주권자가 명령을 내리고, 수범자들은 이에 습관적으로
복종한다''는 명확한 사회 구도는 사라지고, 다수에 의해 내려진 명령에
다수가 복종하거나, 전체 구성원에 의해 내려진 명령에 전체가 복종한다는
모호한 이미지로 대체된다. 그러나 이것은 더 이상 원래 의미의
`명령'(\emphksans{타인(others)}이 특정 방식으로 행동하길 의도한 표현)도
아니고, `복종'도 아니다.

이러한 비판에 대응하기 위해, 사회 구성원들을 개인으로서의 사적
행위능력(private capacity)에서의 개인들과 \textbf{(p.~76)}
선거권자(electors)나 입법자로서의 공적 행위능력에서의 동일한 인격들로
구분할 수 있을 것이다. 이러한 구분은 완전히 이해 가능한 것이며, 실제로
많은 법적·정치적 현상들은 이러한 용어로 가장 자연스럽게 설명된다. 그러나
설령 우리가 한 걸음 더 나아가 공적 행위능력에서의 개인들이 \emphksans{또 다른
인격(another person)}을 구성하며, 그 인격이 습관적으로 복종을 받는다고
말할 준비가 되어 있다 하더라도, 이러한 구분만으로는 주권 이론을 구출할
수 없다. 왜냐하면 어떤 사람들의 집단이 대표자를 선출하거나 명령을 발할
때, 그들이 `개인으로서'가 아니라 `공적 행위능력에서(in their official
capacity)' 행위했다고 말하는 것이 무엇을 의미하는지를 묻는다면, 그에
대한 답변은 오직 특정한 규칙들에 따른 그들의
자격요건들(qualifications)과, 유효한 선거(valid election)나 법을
성립시키기 위해 그들이 따라야 하는 것을 규정하는 다른 규칙들에 대한
준수라는 용어로만 주어질 수 있기 때문이다. 바로 이러한 규칙들에 대한
참조(reference)를 통해서만, 우리는 어떤 행위를 이 사람들의 집단에 의해
이루어진 선거나 법 제정으로 식별할(identify) 수 있다. 이러한 것들은,
개인의 발언이나 서면 명령을 그 개인에게 귀속시키는 데 사용하는 것과
동일한 단순하고 자연적인 기준(simple natural test)에 의해, 그것을
`만든(making)' 집단(body)에게 귀속되는 것이 아니다.

그렇다면 그러한 규칙들이 존재한다는 것은 무엇을 의미하는가? 그것들은
사회 구성원들이 유권자 집단(electorate)으로서 기능하기 위해---그리고
따라서 그 이론의 목적상 주권자로서 기능하기 위해---무엇을 해야 하는지를
규정하는 규칙들이므로, 그 자체가 주권자에 의해 발령된 명령의 지위를 가질
수는 없다. 왜냐하면 그 규칙들이 이미 존재하고 준수되어 있지 않는 한,
어떠한 것도 주권자에 의해 발령된 명령으로 간주될 수 없기 때문이다.

그렇다면 우리는 이러한 규칙들이 인구(population)의 복종의
\emphksans{습관(habits)}에 대한 기술(description)의 일부에 불과하다고 말할 수
있는가? 사회의 대다수가 주권자인 단일한 개인을, 그가 특정한
형식---예컨대 서면으로 작성되고 서명 및 증인이 갖추어진 형식---으로
명령을 내릴 경우에, 그리고 그 경우에만(if, and only if) 복종하는 단순한
사례에서는, (제1절에서 이 맥락에서의 `습관' 개념 사용에 대해 제기된
이의들을 유보한다면) 그가 이러한 방식으로 입법해야 한다는 규칙은 사회의
복종의 습관에 대한 기술의 일부에 불과하다고 말할 수 있을 것이다. 즉,
사회는 그가 이러한 방식으로 명령을 내릴 \emphksans{때(when)} 습관적으로
그에게 복종한다는 것이다. 그러나 주권자인 개인이 규칙들과 독립적으로
식별될 수 없는 경우에는, 이러한 규칙들을 사회가 주권자에게 습관적으로
복종하는 조건이나 항(term)으로서 단지 그렇게 표상할 수는 없다. 이
규칙들은 주권자를 \emphksans{구성하는(constitutive)} 것이지, 주권자에 대한
복종의 습관을 기술할 때 \textbf{(p.~77)} 언급해야 할 사항들에 불과한
것이 아니다. 따라서 현재의 경우, 유권자 집단(electorate)의 절차를
규정하는 규칙들이, 사회가 다수의 개인으로서 유권자 집단으로서의 자기
자신에게 복종하는 조건들을 나타낸다고 말할 수는 없다. 왜냐하면 `유권자
집단으로서의 자기 자신(itself as an electorate)'이라는 표현은 규칙들과
별도로 식별 가능한 어떤 인격(person)에 대한 참조(reference)가 아니기
때문이다. 그것은 선거권자들이 대표자를 선출함에 있어 규칙들을
준수하였다는 사실에 대한 압축된 참조(condensed reference)에 불과하다.
기껏해야 우리는(제1절의 이의들을 유보한다면) 이러한 규칙들이
\emphksans{선출된 인격들(elected persons)}이 습관적으로 복종을 받는 조건들을
제시한다고 말할 수 있을 뿐이다. 그러나 그렇게 되면 우리는
입법부(legislature)가 아니라 유권자 집단이 주권자라는 현재의 설명에서
벗어나, 입법부가 주권자인 이론의 한 형태로 되돌아가게 되며, 그러한
입법부가 그 입법 권력에 대해 법적 제한(legal limitations)을 받을 수
있다는 사실에서 비롯되는 모든 난점들은 여전히 해결되지 않은 채로 남게 될
것이다.

이 장의 앞선 절에서 제시된 논증들과 마찬가지로, 이 이론에 대한 이러한
논증들은, 단지 세부에서 오류가 있다는 주장에 그치지 않고, 명령, 습관,
그리고 복종이라는 단순한 관념만으로는 법에 대한 분석에 충분할 수 없다는
주장을 구성한다는 의미에서 근본적이다. 이에 대신하여 요구되는 것은,
특정한 절차를 준수함으로써 입법할 자격을 갖춘 사람들에게, 제한될 수도
있고 제한되지 않을 수도 있는, 권한을 부여하는 어떤 규칙(a rule
conferring powers)이라는 생각이다.

이 이론이 지니는 일반적 개념적 부적합성과는 별도로, 통상적으로 최고
입법부로 간주되는 기관이 법적으로 제한될 수 있다는 사실을 이 이론 안에
수용하려는 시도에는 많은 보조적 반론들(ancillary objections)이 존재한다.
그러한 경우 주권자를 유권자 집단(electorate)과 동일시해야 한다면, 설령
유권자 집단이 통상 입법부에 대한 모든 제한을 제거할 수 있는 무제한의
개정 권한(unlimited amending power)을 가지고 있는 경우라 하더라도,
그러한 제한들이 법적인 이유가 유권자 집단이 통상 입법부가 습관적으로
복종하는 명령을 부과했기 때문이라고 말할 수 있는지 우리는 물어볼 수
있다. 우리는 입법 권력에 대한 법적 제한들이 명령으로, 그리고 그에 따른
\emphksans{책무들(duties)}로 잘못 표상된다는 우리의 반론을 일단 유보할 수도
있을 것이다. 그렇다 하더라도, 이러한 제한들이 유권자 집단이 입법부에게
이행하라고 심지어 \emphksans{묵시적으로(tacitly)} 명령한 의무라고 가정할 수
있을까? 앞선 장들에서 묵시적 명령(tacit orders)이라는 관념에 대해 제기된
모든 반론들은, 이 맥락에서 그 사용에 대해 더욱 강력한 힘으로 적용된다.
미합중국 헌법에 규정된 것처럼 \textbf{(p.~78)} 그 행사 방식이 극히
복잡한 개정 권한(amending power)을 행사하지 않는다는 사실은, 유권자
집단의 바람들(wishes)을 보여주는 징후로서는 빈약할 수 있으며, 오히려
그들의 무지와 무관심을 보여주는 신뢰할 만한 징후인 경우가 많다. 우리는,
장군이 병사들에게 하사관이 지시하는 바를 따르라고 묵시적으로 명령했다고,
아마도 그럴듯하게, 간주될 수 있는 경우와는 실로 거리가 멀다.

다시 말해, 유권자 집단에 위임된 개정 권한의 범위 바깥에 전적으로 놓여
있는 입법부에 대한 제약이 존재하는 경우, 이 이론의 용어로 우리는
무엇이라고 말해야 하는가? 이는 단지 상상 가능한 경우에 그치지 않고,
실제로 일부 사례들에서는 그러한 상태가 존재한다. 이러한 경우 유권자
집단은 법적 제한의 적용을 받으며, 비록 그것을 통상을 넘어서는
입법부(extraordinary legislature)라고 부를 수는 있을지라도, 법적
제한으로부터 자유롭지 않으므로 주권자는 아니다. 그렇다면 우리는, 사회가
이러한 제한들에 대해 혁명을 일으키지 않았다는 이유로, 전체로서의
사회(the society as a whole)가 주권자이며 이러한 법적 제한들이
\emphksans{그것(it)}에 의해 묵시적으로 명령되어 왔다고 말해야 하는가? 이것이
혁명과 입법 사이의 구별을 유지할 수 없게 만든다는 점만으로도, 아마 이
주장을 배척하기에 충분한 이유가 될 것이다. 마지막으로, 유권자 집단을
주권자로 취급하는 이 이론은, 기껏해야 유권자 집단이 존재하는
민주정에서의 제한된 입법부에 대해서만 설명을 제공한다. 그러나 Rex와 같은
세습 군주(hereditary monarch)가, 체계 내에서 최고(supreme)이면서도
동시에 제한된(limited) 입법 권력을 향유한다는 개념에는 아무런 부조리도
존재하지 않는다.

\subsection{CHAPTER IV 주석}\label{chapter-iv-uxc8fcuxc11d}

50쪽. 오스틴의 주권 이론. 본 장에서 검토된 주권 이론은 오스틴이
『법리학의 범위』(\emph{The Province}) 제5강과 제6강에서 전개한 것이다.
우리는 오스틴이 단지 일정한 형식적 정의들이나 법체계를 논리적으로
배열하기 위한 추상적 도식을 제시한 것이 아니라, 영국이나 미합중국과 같은
법이 존재하는 모든 사회에서는 오스틴이 정의한 속성을 갖춘 주권자가
어딘가에 반드시 존재한다는 \emphksans{사실 주장(factual claim)}을 하고 있다고
해석했다. 다만 이러한 주권자의 존재는 다양한 헌법적·법적 형식들에 의해
가려질 수 있다. 일부 이론가들은 오스틴을 달리 해석하여 그가 그러한 사실
주장을 한 것이 아니라고 본다(예: 스톤, 『법의 영역과 기능』, 제2장 및
제6장, 특히 60, 61, 138, 155쪽. 여기서 오스틴이 여러 공동체에서 주권자를
식별하려 한 시도는 그의 본래 목적에서 벗어난 무의미한 일탈로 다뤄진다).
이러한 해석에 대한 비판은 Morison, `Some Myth about Positivism', loc.
cit., pp.~217--22 참조. 또한 Sidgwick, \emph{The Elements of Politics,}
Appendix (A) `On Austin's Theory of Sovereignty'도 참조.

54쪽. 오스틴에서의 입법적 권위의 연속성. 『\emphksans{법리학의 범위(The
Province)}』에서 `승계의 방식으로 주권을 취득하는 인격들(persons who
``take the sovereignty in the way of succession'')'에 대한 간략한
언급들(제5강(Lecture V), pp.~152--4)은 시사적이지만 불분명하다.
오스틴은, 주권을 취득하는 인물들이 교체되는 일련의 승계를 통해 주권의
연속성을 설명하기 위해서는 그의 핵심 개념들인 `습관적 복종'과
`사령(commands)'에 더하여 무엇인가 추가적인 요소가 필요하다는 점을
인정하는 듯 보이지만, 그 추가 요소를 결코 명확히 식별하지는 않는다. 그는
이와 관련하여 `\emphksans{자격(title)}', 승계에 대한 `\emphksans{청구(claims)}',
그리고 `\emphksans{정당한(legitimate)} 자격'에 관해 말하지만, 이러한 표현들은
통상적으로 사용될 때 모두 승계를 규율하는 \emphksans{규칙}의 존재를 함축하며,
연속적인 주권자들에 대한 단순한 복종의 \emphksans{습관}만을 의미하지는
않는다. 오스틴이 사용하는 이러한 용어들과 `일반적 자격(generic title)'
및 주권을 취득하는 `일반적 방식(the generic mode)'이라는 표현들에 대한
그의 설명은, 주권자의 `확정적(determinate)' 성격에 관한 그의 교설로부터
도출되어야 한다(op. cit., 제5강(Lecture V), pp.~145--55). 여기서 그는
주권자인 개인 또는 개인들이 개별적으로---예컨대 이름으로---식별되는
경우와, `어떤 일반적 기술(generic description)에 부합하는 것으로서'
식별되는 경우를 구별한다. 따라서 (가장 단순한 예를 들면) 세습
군주제에서는 일반적 기술이 특정 조상의 `생존한 장남 후손'일 수 있고,
의회 민주정(parliamentary democracy)에서는 입법부 구성원
자격요건들(qualifications for membership of the legislature)을 재현하는
고도로 복잡한 기술이 될 것이다.

오스틴(Austin)의 견해는, 어떤 사람이 그러한 `일반적' 기술(`generic'
description)을 충족할 때 그는 승계에 대한 `자격(title)' 또는
`권리(right)'를 가진다는 것인 듯하다. 그러나 주권자에 대한 이러한 일반적
\emphksans{기술(description)}에 의존한 설명은, 오스틴이 이 맥락에서
`기술'이라는 말로 승계를 규율하는 수용된 \emphksans{규칙}을 의미하지 않는 한,
그 자체로는 불충분하다. 왜냐하면, 사회의 구성원들이 \emphksans{사실상(as a
matter of fact)} 그때그때 일정한 기술에 부합하는 누구에게든 습관적으로
복종하는 경우와, 그러한 기술에 부합하는 누구든지 복종받을 \emphksans{권리}
또는 \emphksans{자격}을 가진다는 규칙이 수용되어 있는 경우 사이에는 분명한
구별이 존재하기 때문이다. 이는 사람들이 체스 말을 습관적으로 일정한
방식으로 움직이는 경우와, 그렇게 행동할 뿐만 아니라 그것이 말을 움직이는
\emphksans{올바른(right)} 방식이라는 규칙을 수용하고 있는 경우 사이의 차이와
평행(parallel)을 이룬다. 승계에 대한 `권리'나 `자격'이 존재하려면,
승계를 규정하는 규칙(rule providing for the succession)이 반드시 있어야
한다. 오스틴의 일반적 기술(generic descriptions)에 관한 교설은 그러한
규칙을 대체할 수는 없으며, 다만 그 필요성(necessity)을 분명히 드러낼
뿐이다. 입법자로서의 자격을 부여하는 규칙(rule qualifying persons as
legislators)이라는 개념을 인정하지 않은 오스틴에 대한 유사한 비판으로는,
Gray, \emph{The Nature and Sources of the Law}, 제3장(chap.~iii), 특히
제151--157절(ss. 151--7)을 보라. 또한 주권적 기구(sovereign body)의
통일성(unity)과 법인격적(corporate) 또는 `합의체적(collegiate)'
행위능력(capacity)에 관한 오스틴의 제5강(Lecture V)의 설명 역시 동일한
결함을 지니고 있다(이 장 제4절(s. 4) 참조).

55쪽. 규칙과 습관. 여기서 강조된 규칙의 \emphksans{내면적 측면(internal
aspect)}은 제5장 §2(p.88), §3(p.98), 제6장 §1, 제7장 §3에서 더 논의된다.
또한 Hart, ``Theory and Definition in Jurisprudence'', \emph{29 PAS}
Supplementary vol.~(1955), 247--250쪽도 참고. 유사한 관점으로는 Winch,
『사회과학의 개념』(1958), 제2장 57--65쪽, 제3장 84--94쪽, 그리고
Piddington, ``Malinowski의 욕구 이론'', 『인간과 문화』(Firth 편) 참조.

60쪽. 헌법의 기본 규칙에 대한 일반적 수용. 헌법의 수용, 나아가 법체계의
존재와 관련하여, 공무원과 민간 시민들이 법 규칙에 대해 보이는 다양한
태도들은 제5장 §2 (88--91쪽), 제6장 §2 (114--117쪽)에서 더 깊이
논의된다. 또한 Jennings, 『헌법법』(3판) 부록 3: ``법 이론에 대한 노트''
참조.

63쪽. 홉스와 묵시적 사령 이론. 앞선 제3장 §3 및 해당 주석 참조. 또한
Sidgwick, 『정치의 요소들』 부록 A. 현대 입법기관의 제정법조차도
집행되기 전까지는 법이 아니라고 보는 유사한 `현실주의적' 이론으로는
Gray, 『법의 본성과 원천』 제4장, J. Frank, 『법과 현대정신』 제13장
참조.

66쪽. 입법 권한에 대한 법적 제한. 오스틴과 달리, 벤담은 최고 권한이
`명시적 합의(express convention)'에 의해 제한될 수 있으며, 그 합의를
위반하여 제정된 법률은 무효(void)가 될 수 있다고 보았다. 이에 관해서는
『\emphksans{정부에 관한 단편(A Fragment on Government)}』 제4장 제26,
34--38절 참조. 주권자의 권력에 대한 법적 제한의 가능성을 부정하는
오스틴의 논증은, 그러한 제한의 적용을 받는다는 것이 곧
\emphksans{책무(duty)}의 적용을 받는다는 전제에 기초하고 있다. 이에 대해서는
『\emphksans{법리학의 범위(The Province)}』 제6강, 254--268쪽 참조. 그러나
실제로 입법적 권위(legislative authority)에 대한 제한은 책무들(duties)이
아니라 \emphksans{무능력(disabilities)}으로 구성된다(호펠드(Hohfeld),
『\emph{Fundamental Legal Conceptions}』(1923), 제1장 참조).

68쪽. 입법의 방식 및 형식에 관한 규정. 이러한 규정들을 입법 권한에 대한
실질적 제한과 구별하는 것이 어렵다는 점은 제7장 제4절, 149--152쪽에서
추가로 논의된다. 주권 기관의 역량을 `정의하는 것'과 `제한하는 것' 사이의
구분에 대한 포괄적 논의는 마셜, 『의회 주권과
영연방』(\emph{Parliamentary Sovereignty and the Commonwealth}, 1957),
제1--6장 참조.

72쪽. 헌법적 안전장치와 사법 심사. 사법 심사가 허용되지 않는 헌법에
대해서는 휘어, 『현대 헌법』(\emph{Modern Constitutions}) 제7장 참조.
여기에 해당하는 국가는 (주법을 제외한) 스위스, 제3공화국 시기의 프랑스,
네덜란드, 스웨덴 등이 포함된다. 미합중국 대법원이 `정치적 문제(political
questions)'를 야기하는 위헌 주장에 대해 판결을 거부한 사례로는
\emphksans{루터 대 보든}(\emph{Luther v. Borden}), 7 Howard 1, 12 L. Ed. 581
(1849) 참조. 프랑크퍼터의 「대법원」, 『사회과학 백과사전』 제14권,
474--476쪽도 참조.

74쪽. `예외적 입법기관(extraordinary legislature)'으로서의 유권자. 많은
체계에서 통상적 입법기관이 법적 제한을 받는다는 비판을 피하려는 시도로
오스틴이 이 개념을 사용한 것에 대해서는 『법학의 영역』 제6강,
222--233쪽 및 245--251쪽 참조.

76쪽. 사적 행위능력의 입법자와 공적 행위능력의 입법자. 오스틴은 종종
주권 기관의 구성원을 `개별적으로 고려했을 때'와 `구성원으로서 또는
법인적·주권적 행위능력으로 고려했을 때'를 구분한다(『법학의 영역』
제6강, 261--266쪽). 그러나 이 구분은 주권 기관의 입법 활동을 규율하는
규칙 개념을 수반한다. 오스틴은 이러한 공적 또는 합의체적
행위능력(official or collegiate capacity)이라는 생각을 `일반적
기술(generic description)'이라는 불충분한 용어로 설명하려고 암시만 할
뿐이다(54쪽 주석 참조).

78쪽. 개헌 권한의 제한 범위. 미합중국 헌법 제5조의 단서 조항 참조. 독일
연방공화국의 기본법(1949년) 제1조와 제20조는 제79조 제3항이 부여한 개헌
권한의 범위 밖에 놓여 있다. 터키 헌법(1945년)의 제1조 및 제102조도 참조.

\subsection{CHAPTER IV 제3판
주석}\label{chapter-iv-uxc81c3uxd310-uxc8fcuxc11d}

55--57쪽. \emphksans{사회 규칙의 본성에 대하여}. 하트의 `실천 이론(practice
theory)'에 대한 비판은 다음 저작들을 참조할 것: G. J. 워녹(G. J.
Warnock), 『도덕성의 대상(The Object of Morality)』(메튜언, 1971) 제4장;
로널드 드워킨(Ronald Dworkin), 『권리를 진지하게 대하기(Taking Rights
Seriously)』 48--58쪽; 조셉 라즈(Joseph Raz), 『실천적 이성과
규범(Practical Reason and Norms)』 49--58쪽.

하트는 이후의 저술에서 사령에 관한 사고를 발전시켜 규칙에 대한 또 다른
견해를 제안하였다. 그는 사령이 `결정적인(peremptory)' 것으로서, 즉
``청자가 어떤 행위를 할 것인지에 대해 그 장단점을 독자적으로 숙고하는
것을 차단하거나 차단하려는 의도를 가진 것''이며, 또한 `내용
독립적(content-independent)'으로, ``실행될 행위의 성격이나 본질에
관계없이 이유로 기능하도록 의도된 것''이라고 말한다. 그는 이어,
``사회에서 사령자의 발언이 결정적인 행위 이유(peremptory reasons for
action)로 일반적으로 승인되는 것은 사회 규칙의 존재와 동일하다''고
주장한다(H. L. A. Hart, 「사령과 권위 있는 법적 이유(Commands and
Authoritative Legal Reasons)」, 『벤담에 관한 에세이( \emph{Essays on
Bentham})』 제10장, 258쪽 참조). 하트는 `법적·도덕적 의무(Legal and
Moral Obligation)'라는 논문에서 `내용 독립적 이유' 개념을 처음
도입하였으며, 이는 A. I. 멜든(A. I. Melden) 편 『도덕철학 에세이(Essays
in Moral Philosophy)』(워싱턴 대학 출판부, 1958), 82--107쪽에 수록되어
있다. 이 개념은 라즈의 『자유의 도덕성(The Morality of
Freedom)』(옥스퍼드 대학 출판부, 1986) 제2장에서, 그리고 레슬리
그린(Leslie Green)의 『국가의 권위(The Authority of the
State)』(옥스퍼드 대학 출판부, 1990), 36--42쪽에서 더욱 발전되었으며, P.
마크위크(P. Markwick)의 「법과 내용 독립적 이유(Law and
Content-Independent Reasons)」(\emph{Oxford Journal of Legal Studies},
2000년 제20권, 579쪽)에서는 비판받았다.

56--58쪽. \emphksans{법의 `내면적 측면'(internal aspect)}. 하트는 종종
`내면적 관점(internal point of view)'과 `내면적 측면(internal aspect)'을
혼용해서 사용한다(예: 88--89쪽 참조). 이 일반 개념에 대해서는 닐
매코믹(Neil MacCormick), 『H. L. A. Hart』 29--40쪽 참조. 이 개념의
모호성에 대한 논의로는 스티븐 페리(Stephen Perry), 「하트의 사회 규칙과
법의 기초: 내면적 관점의 해방(Hart on Social Rules and the Foundations
of Law: Liberating the Internal Point of View)」(\emph{Fordham Law
Review}, 2006년 제75권, 1171쪽), 스콧 샤피로(Scott J. Shapiro), 「내면적
관점이란 무엇인가?(What is the Internal Point of View?)」(\emph{Fordham
Law Review}, 2006년 제75권, 1157쪽) 참조. 법이론에서 이 개념의 방법론적
중요성은 라즈, 『권위와 해석 사이(Between Authority and
Interpretation)』 제2장에서 논의된다. 이 방법론적 주장에 대한 비판으로는
스티븐 페리, 「법이론에서의 해석과 방법론(Interpretation and Methodology
in Legal Theory)」, 안드레이 마머(Andrei Marmor) 편 『법과 해석(Law and
Interpretation)』(옥스퍼드 대학 출판부, 1997), 존 피니스(John Finnis),
『자연법과 자연권(Natural Law and Natural Rights)』 제1장 등이 있다.
피니스의 비판은 줄리 딕슨(Julie Dickson), 『평가와 법이론(Evaluation and
Legal Theory)』(하트 출판사, 2001), 제2--4장에서 면밀히 분석된다.

66--71쪽. \emphksans{입법권에 대한 법적 제한}. 하트는 『벤담에 관한 에세이(
\emph{Essays on Bentham})』 제9장 「주권과 법적으로 제한된
정부(Sovereignty and Legally Limited Government)」에서 벤담이 입법
제한을 이해하려 한 시도를 분석한다. 오스틴의 무제한성(ilimitability)
개념에 대한 라즈의 비판은 『법체계의 개념(The Concept of a Legal
System)』 제2장에서 다루어진다. 오스틴의 분석이 헌법이론에 미치는 함의에
대해서는 제프리 마셜(Geoffrey Marshall), 『헌법 이론(Constitutional
Theory)』(옥스퍼드 대학 출판부, 1980), 제1, 3장 참조.
